% !TeX root = ../../Book2.tex
\chapter{Strichartz 估计与色散型偏微分方程}
\label{b2:app:E}

\begin{chapterabstract}
Schrödinger 演化保持平方范数，却能使局部集中的振幅随时间散开。色散界描述这种振荡效应，Strichartz 估计则把空间可积性与时间可积性结合起来。本附录从 Gaussian 正则化推导核公式，以 \(TT^*\) 和时间分数积分证明非端点估计，并建立推迟积分的连续延拓。半波演化还需先作频率局部化，再以平方函数合成；能量端点所需的深层双线性求和另行标源。最后用已证的非端点工具构造一维三次方程的温和解，通过保质量的频率截断把守恒律传给平方可积初值，完成全局延拓。
\end{chapterabstract}

\begin{learninggoals}\begin{itemize}
\item\label{b2:item:app-E-goal-1}从正则化核公式取得色散界，区分固定时间范数与时间空间混合范数。
\item\label{b2:item:app-E-goal-2}完成非端点 \(TT^*\) 对偶论证，并在适当外力空间中定义推迟积分。
\item\label{b2:item:app-E-goal-3}由环带驻相、频率尺度和平方函数合成推出明确范围的半波估计。
\item\label{b2:item:app-E-goal-4}构造一维三次方程的局部解，利用频率截断证明质量守恒、全局延拓与连续依赖。
\end{itemize}\end{learninggoals}

所需基础是本卷的 Plancherel 定理与复插值、\chapref{b2:ch:19}的分数积分和平方函数，以及\chapref{b2:ch:21}的复 Gaussian 积分与驻相法。\appref{b2:app:D}解释了酉群、生成元和温和解的关系；本附录仍会明确给出 Fourier 乘子和积分意义，避免把无界算子的形式求导当作粗糙初值的证明。空间函数取复值，测度均为相应 Euclidean 空间的 Lebesgue 测度，Fourier 约定保持为 \(e^{-2\pi\mathrm i x\cdot\xi}\)。

\ProseParagraphBreak

前两节先建立 Schrödinger 的线性估计，第三节说明波方程为何多出频率权重，第四节区分本附录已经证明的非端点范围与另需深层理论的端点。最后一节的非线性应用只依赖前两节的完整证明，因此可以先读这条路线，再返回波估计和端点。局部收缩、质量守恒和全局解各有不同的论证任务：时间空间估计提供收缩所需的范数，逼近使守恒律对一般初值成立，统一的局部存在时间才允许反复延拓。

\section{Schrödinger 色散与混合范数}
\label{b2:sec:app-E-1}

热演化通过频率乘子的模衰减压低高频；自由 Schrödinger 演化的乘子则始终具有单位模。它保存平方能量，却仍可能使一个局部波包的峰值下降。这里要估计的正是振荡相位造成的空间分散，再把这种固定时间的估计与时间积分结合起来。先从 Fourier 定义出发，可以避免把不绝对可积的振荡核误当成热核使用。

\Needspace{6\baselineskip}
\begin{proposition}[酉群与色散核]\label{b2:prop:E-dispersive-kernel}
设整数 \(d\geq1\)，所有函数均为复值，空间测度为 \(\mathbb R^d\) 上的 Lebesgue 测度。采用
\[
 \widehat f(\xi)=\int_{\mathbb R^d}
       e^{-2\pi\mathrm ix\cdot\xi}\cdot f(x)\,\mathrm dx
\]
及其 Plancherel 延拓 \(\mathcal F_2\)。对 \(f\in L^2(\mathbb R^d)\) 定义
\[
 U(t)f=\mathcal F_2^{-1}
       \left(e^{-4\pi^2\mathrm it|\xi|^2}\mathcal F_2f\right),
 \quad t\in\mathbb R.
\]
则 \(U\) 为强连续酉群，且
\[
 U(t+s)=U(t)U(s),\quad U(t)^*=U(-t),\quad
 \|U(t)f\|_2=\|f\|_2.
\]
对固定 \(t\ne0\)，令
\begin{equation}\label{b2:eq:E-fresnel-kernel}
 \begin{aligned}
 K_t(x)&=(4\pi\mathrm it)^{-d/2}
               \exp\left(\frac{\mathrm i|x|^2}{4t}\right),\\
 (4\pi\mathrm it)^{-d/2}
 &=(4\pi|t|)^{-d/2}
       \exp\left(-\frac{\mathrm i\pi d}{4}\operatorname{sgn}t\right).
 \end{aligned}
\end{equation}
对 \(f\in L^1\cap L^2\)，有 \(U(t)f=K_t*f\) 几乎处处。卷积公式还定义从 \(L^1\) 到 \(L^\infty\) 的有界算子，满足
\[
 \|U(t)f\|_\infty\leq(4\pi|t|)^{-d/2}\|f\|_1.
\]
它与 \(L^2\) 上的酉群在共同定义域相容。对 \(2\leq r\leq\infty\)，令 \(r'\) 为共轭指数，则相容插值还给出
\begin{equation}\label{b2:eq:E-dispersive-interpolation}
 \|U(t)f\|_r\leq C_{d,r}|t|^{-\alpha}\|f\|_{r'},
 \quad
 \alpha=d\left(\frac12-\frac1r\right),\quad t\ne0.
\end{equation}
\end{proposition}
\begin{proof}
乘子的模为一，Plancherel 定理给出等距性；乘子相乘和取共轭分别给出群律与伴随公式。特别地，\(U(-t)\) 是逆，所以算子是酉的。对 \(h\to0\)，
\[
 \|(U(t+h)-U(t))f\|_2^2
 =\int_{\mathbb R^d}
   |e^{-4\pi^2\mathrm ih|\xi|^2}-1|^2
        \cdot|\mathcal F_2f(\xi)|^2\,\mathrm d\xi.
\]
被积函数逐点趋零，并由 \(4|\mathcal F_2f|^2\) 控制，支配收敛给出全部时间的强连续性。

为取得核，固定 \(\varepsilon>0\)，把乘子改成
\[
 m_{\varepsilon,t}(\xi)
 =e^{-4\pi^2(\varepsilon+\mathrm it)|\xi|^2}.
\]
它在频率空间可积。逐坐标使用\cref{b2:prop:21-complex-gaussian}，参数取 \(a=4\pi(\varepsilon+\mathrm it)\)，频率取 \(-x_j\)。该命题要求 \(\operatorname{Re}a>0\)，这里由 \(\varepsilon>0\) 保证；平方根在右半平面选择与正实轴正平方根连续相接的分支。因此逆变换为
\begin{equation}\label{b2:eq:E-regularized-kernel}
 K_{\varepsilon,t}(x)
 =[4\pi(\varepsilon+\mathrm it)]^{-d/2}
       \exp\left(-\frac{|x|^2}{4(\varepsilon+\mathrm it)}\right).
\end{equation}
对每个 \(\varepsilon>0\)，指数实部为负的二次型，所以 \(K_{\varepsilon,t}\in L^1\)，卷积与 Fourier 乘子定义相容。先在 Schwartz 函数上由绝对积分验证，再由 Young 不等式、Plancherel 和稠密性，便得到 \(L^2\) 输入上的相容性。

现在固定 \(t\ne0\)，令 \(\varepsilon\downarrow0\)。复数 \(\varepsilon+\mathrm it\) 的辐角趋于 \((\pi/2)\operatorname{sgn}t\)，所以复幂趋于\eqref{b2:eq:E-fresnel-kernel}规定的分支；指数趋于 \(\exp(\mathrm i|x|^2/(4t))\)。另外，
\[
 |K_{\varepsilon,t}(x)|
 =[4\pi|\varepsilon+\mathrm it|]^{-d/2}
       \exp\left(-\frac{\varepsilon|x|^2}
                       {4(\varepsilon^2+t^2)}\right)
 \leq(4\pi|t|)^{-d/2}.
\]
于是对 \(f\in L^1\cap L^2\) 和每个固定 \(x\)，卷积积分由
\((4\pi|t|)^{-d/2}|f(y)|\) 控制，故
\[
 (K_{\varepsilon,t}*f)(x)\longrightarrow(K_t*f)(x).
\]
另一方面，频率侧乘子逐点趋于 \(e^{-4\pi^2\mathrm it|\xi|^2}\)，差平方由 \(4|\widehat f|^2\) 控制。因此相同的正则化算子在 \(L^2\) 中趋于 \(U(t)f\)。选一个几乎处处收敛的子列，便把逐点卷积极限与酉群极限识别为同一函数。

对一般 \(L^1\) 输入，\(K_t\) 的模是常数，故 \(\int K_t(x-y)\cdot f(y)\,\mathrm dy\) 对每个 \(x\) 绝对收敛，并给出所述 \(L^\infty\) 界。用 \(L^1\cap L^2\) 输入在 \(L^1\) 中逼近，也得到同一个延拓。最后在这个 \(L^1\to L^\infty\) 算子与 \(L^2\to L^2\) 酉算子之间应用 Riesz--Thorin 插值。取插值参数 \(2/r\)，则输入指数为 \(r'\)，输出指数为 \(r\)，时间幂为
\[
 \frac d2\left(1-\frac2r\right)
 =d\left(\frac12-\frac1r\right).
\]
这就得到\eqref{b2:eq:E-dispersive-interpolation}。两个端点算子在交空间上一致，正是这一步插值所需的相容性。
\end{proof}
核 \(K_t\) 自身在全空间的绝对积分发散，但它与一个 \(L^1\) 输入的卷积仍绝对收敛。这两件事并不矛盾。群律来自 Fourier 乘子，不能未经核验就把两个非可积振荡核作普通绝对卷积来证明它。以上论证也先固定了 \(t\ne0\)，再去掉 Gaussian 正则化；所得控制在 \(t\to0\) 时并不一致。

\ProseParagraphBreak

固定时间的色散界把较强的输入可积性换成输出的空间范数。它本身还不能处理只有 \(L^2\) 的初值：一般 \(L^2\) 函数不属于 \(L^{r'}\)，而非线性方程中的初值常只能由平方能量控制。要取得这类输入的空间可积性，就必须同时衡量它在时间上出现的程度。

\Needspace{6\baselineskip}
\begin{definition}\label{b2:def:E-mixed-norm}
设整数 \(d\geq1\)，\(I\subset\mathbb R\) 为区间，\(u:I\times\mathbb R^d\to\mathbb C\) 为 Lebesgue 可测函数。对 \(1\leq q,r<\infty\)，定义\term[hunhe-fanshu]{混合范数}[mixed norm]
\[
 \|u\|_{L^q_tL^r_x(I\times\mathbb R^d)}
 =\left(\int_I
     \left(\int_{\mathbb R^d}|u(t,x)|^r\,\mathrm dx\right)^{q/r}
                   \mathrm dt\right)^{1/q}.
\]
当 \(r=\infty\) 时，内层范数改为关于 \(x\) 的本性上确界；当 \(q=\infty\) 时，外层范数改为关于 \(t\) 的本性上确界。范数有限者组成 \(L^q_tL^r_x(I\times\mathbb R^d)\)，也记为 \(L^q(I;L^r(\mathbb R^d))\)。
当 \(r=\infty\) 时，此记号按可测函数的迭代范数解释，不额外要求取值于 \(L^\infty\) 的映射具有 Bochner 强可测性；下文需要稠密性与强积分的论证均另行给出适用指数。
\end{definition}
\symidx{\(L^q_tL^r_x\)：先取空间范数、再取时间范数的混合空间}
这里的次序是先在一个固定时刻计算空间大小，再把这条随时间变化的大小曲线放入 \(L^q(I)\)。若 \(u(t,x)=a(t)b(x)\)，则
\[
 \|u\|_{L^q_tL^r_x}=\|a\|_{L^q(I)}\cdot\|b\|_{L^r(\mathbb R^d)}.
\]
当 \(q=r\) 时，Tonelli 定理使它等于通常的时空 \(L^q\) 范数；不同指数时，一般不能直接交换两个范数的次序。非线性乘积也必须分两步估计：先核对每个时间的空间 Hölder 指数，再核对所得时间函数的 Hölder 指数。

\ProseParagraphBreak

例如，\(\|u\|_{L^\infty_tL^2_x}\) 要求几乎每个时刻都有共同的能量界；有限的 \(q\) 则允许个别时段的空间范数很大，只要其 \(q\) 次方在时间上可积。对于连续 \(L^2\) 轨道，时间本性上确界等于通常的上确界，而一般 \(L^q_tL^r_x\) 元素只按几乎处处意义定义，不能自动给每个时刻赋予一个 \(L^r\) 值。

\Needspace{6\baselineskip}
\begin{proposition}[Schrödinger 缩放]\label{b2:prop:E-schrodinger-scaling}
设整数 \(d\geq1\)，\(f\in L^2(\mathbb R^d;\mathbb C)\)，\(U(t)=e^{\mathrm it\Delta}\) 采用\cref{b2:prop:E-dispersive-kernel}的 Fourier 定义。对 \(\lambda>0\)，令
\[
 f_\lambda(x)=\lambda^{d/2}f(\lambda x),\quad u(t)=U(t)f.
\]
则 \(\|f_\lambda\|_2=\|f\|_2\)，并且
\[
 U(t)f_\lambda(x)
 =u_\lambda(t,x)
 =\lambda^{d/2}u(\lambda^2t,\lambda x).
\]
对 \(1\leq q,r\leq\infty\) 和任意时间区间 \(I\)，有
\begin{equation}\label{b2:eq:E-schrodinger-scaling}
 \|u_\lambda\|_{L^q_tL^r_x(I)}
 =\lambda^{d/2-d/r-2/q}
       \|u\|_{L^q_tL^r_x(\lambda^2I)},
\end{equation}
其中 \(1/\infty=0\)。因此，若存在与 \(f\) 无关的全时间估计
\[
 \|U(t)f\|_{L^q_tL^r_x(\mathbb R\times\mathbb R^d)}
 \leq C\|f\|_2
 \quad(f\in L^2),
\]
则必有
\[
 \frac2q+\frac dr=\frac d2.
\]
\end{proposition}
\begin{proof}
空间换元给出初值范数不变，而 Fourier 换元给出
\(\widehat{f_\lambda}(\xi)=\lambda^{-d/2}\widehat f(\xi/\lambda)\)。把 \(\xi=\lambda\eta\) 代入演化乘子，就得到所述二次时间缩放。先对 Schwartz 初值计算，再由酉性和 \(L^2\) 稠密性推广。

固定时间的空间范数产生因子 \(\lambda^{d/2-d/r}\)，随后在时间积分中令 \(s=\lambda^2t\)，再产生因子 \(\lambda^{-2/q}\)。无穷指数直接按本性上确界换元，得到同一公式。

为证明必要性，假设给定的全时间界成立，选一个非零 Gaussian 初值 \(f\)。其输出混合范数由假设有限，并且大于零。把估计应用于每个 \(f_\lambda\)，得到
\[
 \lambda^{d/2-d/r-2/q}\|U(t)f\|_{L^q_tL^r_x}
 \leq C\|f\|_2\quad(\lambda>0).
\]
若幂指数非零，令 \(\lambda\to\infty\) 或 \(\lambda\downarrow0\)，左端就会发散。因此该幂只能为零。
\end{proof}
必要性使用的是全时间区间及与缩放无关的常数。在固定有限区间上，缩放同时把区间变成 \(\lambda^2I\)，所以不能删去区间变化后仍照搬同一论证。指数关系也只筛选候选估计，并不提供充分性证明。

\Needspace{6\baselineskip}
\begin{definition}\label{b2:def:E-admissible}
设整数 \(d\geq1\)。若 \(2\leq q,r\leq\infty\)，
\[
 \frac2q+\frac dr=\frac d2,\quad(q,r,d)\ne(2,\infty,2),
\]
其中 \(1/\infty=0\)，就称 \((q,r)\) 为 \(d\) 维的\term[schrodinger-kerongxu-dui]{Schrödinger 可容许对}[Schrödinger admissible pair]。
\end{definition}
例如一维的 \((8,4)\)、二维的 \((4,4)\) 以及三维的 \((4,3)\) 都满足关系式。能量对 \((\infty,2)\) 在每个维数都可用。下面完整证明的范围是 \(2<q<\infty\)、\(r<\infty\)，另加能量对；可容许对的定义本身并没有把其余边界情形一并证明。

\ProseParagraphBreak

Gaussian 初值可以把衰减与能量保存放在同一个算式中。取 \(f(x)=e^{-|x|^2}\)，复 Gaussian 公式给出
\[
 U(t)f(x)=(1+4\mathrm it)^{-d/2}
        \exp\left(-\frac{|x|^2}{1+4\mathrm it}\right),
\]
故
\[
 |U(t)f(x)|=(1+16t^2)^{-d/4}
        \exp\left(-\frac{|x|^2}{1+16t^2}\right).
\]
峰高下降，典型宽度却按 \((1+16t^2)^{1/2}\) 增长。对 \(2\leq r<\infty\) 直接作空间换元，得到
\[
 \|U(t)f\|_r
 =\left(\frac\pi r\right)^{d/(2r)}
        (1+16t^2)^{-\alpha/2},
 \quad \alpha=d\left(\frac12-\frac1r\right).
\]
\(r=2\) 时，时间幂为零，范数保持为 \((\pi/2)^{d/4}\)；\(r=\infty\) 时直接取峰值得到相应极限公式。对有限 \(q\)，大时间的 \(q\) 次方可积要求 \(q\alpha>1\)，而可容许关系给出 \(q\alpha=2\)。这个算例验证了候选估计与传播尺度相容，尚不能代替对所有 \(L^2\) 初值的证明。

\ProseParagraphBreak

最后，任何固定时间的酉性都不会自行给出 \(L^2\to L^r\)、\(r>2\) 的有界算子。若有这样的界，利用 \(U(t)\) 在 \(L^2\) 上满射，就会得到 \(\|g\|_r\leq C\|g\|_2\) 对所有 \(g\in L^2\) 成立。取一个非零光滑紧支函数，按 \(g_\lambda(x)=\lambda^{d/2}g(\lambda x)\) 缩放，右端保持不变而左端按 \(\lambda^{d(1/2-1/r)}\) 增长，产生矛盾。Strichartz 估计能成立，正因为它保留了时间积分，没有要求每一个固定时刻都具备这种不可能的空间映射界。

\section{非端点 \texorpdfstring{\(TT^*\)}{TT*} 证明}
\label{b2:sec:app-E-2}

现在希望从平方可积初值取得有限指数的时空范数。直接对 \(U(t)f\) 套用色散界会要求 \(f\in L^{r'}\)，与给定初值不符。\(TT^*\) 方法先把时空测试函数送回 \(L^2\)，再把两个时间之间的演化写成 \(U(t-s)\)。这样空间色散恰好作用于测试函数的 \(L^{r'}\) 范数，而时间变量留下一个可由分数积分控制的正核。

\ProseParagraphBreak

本节的复内积取第一变量线性：
\[
 \langle f,g\rangle_{L^2}
 =\int_{\mathbb R^d}f(x)\cdot\overline{g(x)}\,\mathrm dx.
\]
时空配对也按 \(\int F\cdot\overline G\) 排列。有限指数 \(1<q,r<\infty\) 时，混合 Hölder 不等式先在空间、再在时间给出
\[
 \left|\int_{\mathbb R}\int_{\mathbb R^d}
                u\cdot\overline F\,\mathrm dx\,\mathrm dt\right|
 \leq\|u\|_{L^q_tL^r_x}\cdot\|F\|_{L^{q'}_tL^{r'}_x}.
\]
对应的范数对偶公式可以先限制在有限时空区域内，用两次 Hölder 的取等方向得到，再以截断和光滑逼近扩展到全空间。因而控制所有紧支光滑测试函数的配对，可以检测 \(L^q_tL^r_x\) 范数；不必事先假设待估计函数已有有限的全时间范数。

\Needspace{6\baselineskip}
\begin{theorem}[非端点 Schrödinger Strichartz]\label{b2:thm:E-schrodinger-nonendpoint}
设整数 \(d\geq1\)，
\[
 U(t)f=\mathcal F_2^{-1}
       \left(e^{-4\pi^2\mathrm it|\xi|^2}\mathcal F_2f\right),
 \quad f\in L^2(\mathbb R^d;\mathbb C).
\]
若 \(2<q<\infty\)、\(2\leq r<\infty\) 且 \(2/q+d/r=d/2\)，则
\begin{equation}\label{b2:eq:E-homogeneous-strichartz}
 \|U(t)f\|_{L^q_tL^r_x(\mathbb R\times\mathbb R^d)}
 \leq C_{d,q,r}\|f\|_2.
\end{equation}
此外，\((q,r)=(\infty,2)\) 时有能量恒等式。对上述有限指数对，令 \(q',r'\) 为共轭指数；在紧支光滑函数上定义的算子
\[
 T^*F=\int_{\mathbb R}U(-s)F(s)\,\mathrm ds
\]
唯一连续延拓为从 \(L^{q'}(\mathbb R;L^{r'}(\mathbb R^d))\) 到 \(L^2(\mathbb R^d)\) 的算子，且
\begin{equation}\label{b2:eq:E-adjoint-bound}
 \|T^*F\|_2\leq C_{d,q,r}\|F\|_{L^{q'}_tL^{r'}_x}.
\end{equation}
\end{theorem}
\begin{proof}
先只取 \(F\in C_c^\infty(\mathbb R\times\mathbb R^d)\)。其时间支撑紧，且 \(s\mapsto F(s)\) 连续取值于 \(L^2\)，所以 \(T^*F\) 作为 \(L^2\) 值强积分已经有定义；此处没有使用待证的时空估计。记
\[
 V_F(t)=U(t)T^*F
       =\int_{\mathbb R}U(t-s)F(s)\,\mathrm ds.
\]
后一个等式来自有界算子移入强积分与群律。对光滑紧支输入，它也是每个固定时间的 \(L^r\) 积分；可以先在 Schwartz 类逐项计算，再使用下述可积的上界。

令
\[
 \alpha=d\left(\frac12-\frac1r\right)=\frac2q.
\]
由 \(q>2\) 得 \(0<\alpha<1\)。色散插值界\eqref{b2:eq:E-dispersive-interpolation}与 Minkowski 不等式给出
\begin{equation}\label{b2:eq:E-ttstar-positive-kernel}
 \|V_F(t)\|_r
 \leq C\int_{\mathbb R}|t-s|^{-\alpha}
                     \|F(s)\|_{r'}\,\mathrm ds.
\end{equation}
对紧支光滑输入，时间奇点局部可积，右端逐点有意义。把时间看成一维空间，核 \(|t|^{-\alpha}\) 是阶 \(\beta=1-\alpha\) 的分数积分核，差一个固定正常数。\cref{b2:thm:19-hls}适用，因为
\[
 0<\beta<1,\quad
 1<q'<\frac1\beta,\quad
 \frac1q=\frac1{q'}-\beta.
\]
中间的不等式可由 \(\beta q'=(q-2)/(q-1)<1\) 直接核验。因此
\begin{equation}\label{b2:eq:E-ttstar-bound}
 \|V_F\|_{L^q_tL^r_x}
 \leq C\|F\|_{L^{q'}_tL^{r'}_x}.
\end{equation}

下一步才取得 \(T^*\) 的界。按第一变量线性的内积约定，酉性给出
\[
 \begin{aligned}
 \|T^*F\|_2^2
 &=\int_{\mathbb R}
      \langle U(t)T^*F,F(t)\rangle_{L^2}\,\mathrm dt\\
 &\leq \|V_F\|_{L^q_tL^r_x}
              \cdot\|F\|_{L^{q'}_tL^{r'}_x}\\
 &\leq C\|F\|_{L^{q'}_tL^{r'}_x}^2.
 \end{aligned}
\]
第一行也可把 \(U(t)\) 移到第二变量，再由强积分直接核验；时间紧支与 \(L^2\) 有界性保证换序合法。由于 \(q',r'\) 都有限，紧支光滑函数在输入混合空间中稠密，所以这个估计唯一地定义一般 \(F\) 的 \(T^*F\)。一般输入不需要预先属于 \(L^1_tL^2_x\)。

最后取 Schwartz 初值 \(f\)。对上述光滑测试函数，
\[
 \int_{\mathbb R}
       \langle U(t)f,F(t)\rangle_{L^2}\,\mathrm dt
 =\langle f,T^*F\rangle_{L^2}.
\]
右端的绝对值至多为
\(C\|f\|_2\cdot\|F\|_{L^{q'}L^{r'}}\)。应用混合范数对偶公式，先在有限时空区域内估计，再扩大区域，得到\eqref{b2:eq:E-homogeneous-strichartz}，没有预设 \(T:f\mapsto U(t)f\) 的有界性。

对一般 \(f\in L^2\)，取 Schwartz 函数 \(f_n\to f\)。已经得到的界使 \(U(t)f_n\) 在 \(L^qL^r\) 中 Cauchy。酉性又给出
\[
 \sup_{t\in\mathbb R}\|U(t)f_n-U(t)f\|_2
 =\|f_n-f\|_2\longrightarrow0.
\]
两个极限在紧支光滑时空测试函数上的配对相同，因而在分布意义、进而几乎处处意义下相同。所构造的混合空间元素就是原先定义的酉轨道。能量对则直接由酉性成立。
\end{proof}
证明的次序是：先对光滑测试函数定义 \(T^*\)，再以色散和 HLS 控制 \(U(t)T^*F\)，从二次配对取得有界 \(T^*\)，最后对偶得到 \(T\)。符号 \(TT^*\) 描述的是最终算子结构；若一开始便假设 \(T\) 已有目标有界性，就会把待证结论放进前提。

\ProseParagraphBreak

当 \(q=2\) 时，\(\alpha=1\)，正时间核变成 \(|t-s|^{-1}\)，不再处于上述一维 HLS 的阶数范围。这说明当前取绝对值的证明不能达到该端点，并不证明所有端点估计都失败。有限空间指数也是本次对偶证明的一项条件；例如一维可容许对 \((4,\infty)\) 未包含在本定理的有限指数结论中。全时间齐次界限制到任意区间仍成立，且时间平移给出
\[
 \|U(t-a)f\|_{L^q_tL^r_x(I)}\leq C\|f\|_2
 \quad(a\in\mathbb R),
\]
常数不依赖 \(I\) 或起点 \(a\)。

\ProseParagraphBreak

非齐次演化只累计早于观察时刻的外力。对光滑输入可以写出通常的迟滞积分；对一般 \(L^{q'}_tL^{r'}_x\) 输入，其每个时间的空间值未必属于 \(L^2\)，所以不能预先把同一个符号解释为 \(L^2\) 值绝对积分。先使用已经建立的 \(T^*\) 延拓，可以同时定义每个时间的输出并选择连续的 \(L^2\) 代表。

\Needspace{6\baselineskip}
\begin{proposition}[同对非齐次估计]\label{b2:prop:E-retarded}
设整数 \(d\geq1\)，\(U(t)\) 为符号 \(e^{-4\pi^2\mathrm it|\xi|^2}\) 的 Schrödinger 酉群。设 \(2<q<\infty\)、\(2\leq r<\infty\)、\(2/q+d/r=d/2\)，且 \(q',r'\) 为共轭指数。令 \(I\subset\mathbb R\) 为任意区间，\(F\in L^{q'}(I;L^{r'}(\mathbb R^d;\mathbb C))\)，并记 \(\widetilde F\) 为在 \(I\) 外补零的函数。以\cref{b2:thm:E-schrodinger-nonendpoint}中由光滑强积分延拓的算子 \(T^*\)，定义
\[
 B_F(t)=T^*\left(\mathbf1_{(-\infty,t)}\cdot\widetilde F\right),
 \quad \mathcal R_I F(t)=U(t)B_F(t).
\]
则 \(\mathcal R_I F\in C_b(\mathbb R;L^2)\cap L^q(\mathbb R;L^r)\)，并且
\begin{equation}\label{b2:eq:E-retarded-bound}
 \|\mathcal R_I F\|_{L^\infty_tL^2_x(\mathbb R)}
 +\|\mathcal R_I F\|_{L^q_tL^r_x(\mathbb R)}
 \leq C\|F\|_{L^{q'}_tL^{r'}_x(I)}.
\end{equation}
它是光滑输入上迟滞积分
\[
 \mathcal R_I F(t)=\int_{s<t,\ s\in I}
                         U(t-s)F(s)\,\mathrm ds
\]
的连续延拓。对任意有限 \(a\in\mathbb R\)，还可定义有限起点的有向积分
\begin{equation}\label{b2:eq:E-finite-start-duhamel}
 \begin{aligned}
 \mathcal D_aF(t)
 &\coloneqq U(t)(B_F(t)-B_F(a))\\
 &=\mathcal R_I F(t)-U(t-a)\mathcal R_I F(a).
 \end{aligned}
\end{equation}
光滑情形它等于 \(\int_a^tU(t-s)\widetilde F(s)\,\mathrm ds\)，当 \(t<a\) 时按负的反向积分解释。它具有相同的 \(C_bL^2\) 与 \(L^qL^r\) 控制，常数仍只依赖 \(d,q,r\)。特别地，限制到 \(I\) 后，所有界都与区间长度及有限起点无关，且 \(\mathcal D_aF(a)=0\)。
\end{proposition}
\begin{proof}
由\eqref{b2:eq:E-adjoint-bound}及酉性，对每个 \(t\) 有
\[
 \|\mathcal R_I F(t)\|_2
 \leq C\|\mathbf1_{(-\infty,t)}\cdot\widetilde F
                         \|_{L^{q'}L^{r'}}
 \leq C\|F\|_{L^{q'}L^{r'}(I)}.
\]
若 \(h\to0\)，两次截断的差只支撑在 \(t\) 与 \(t+h\) 之间，因而
\[
 \|B_F(t+h)-B_F(t)\|_2
 \leq C\|\mathbf1_{(\min\{t,t+h\},\max\{t,t+h\})}
                       \cdot\widetilde F\|_{L^{q'}L^{r'}}
 \longrightarrow0.
\]
最后的收敛来自 \(\|\widetilde F(s)\|_{r'}^{q'}\) 的积分绝对连续性，而不是外力逐点连续。再写
\[
 \begin{aligned}
 \mathcal R_I F(t+h)-\mathcal R_I F(t)
 &=U(t+h)(B_F(t+h)-B_F(t))\\
 &\quad+(U(t+h)-U(t))B_F(t).
 \end{aligned}
\]
第一项由酉性趋零，第二项是固定 \(L^2\) 向量的群轨道差。因此所得输出在全部实时间连续，包括原区间的有限端点。

为证明混合范数界，先取 \(G\in C_c^\infty(\mathbb R\times\mathbb R^d)\)。普通迟滞积分满足
\[
 \left\|\int_{s<t}U(t-s)G(s)\,\mathrm ds\right\|_r
 \leq C\int_{s<t}|t-s|^{-\alpha}\|G(s)\|_{r'}\,\mathrm ds
 \leq C\int_{\mathbb R}|t-s|^{-\alpha}\|G(s)\|_{r'}\,\mathrm ds.
\]
同一个 HLS 估计因而适用，常数不随截断位置改变。现在选光滑紧支 \(G_n\to\widetilde F\) 于 \(L^{q'}L^{r'}\)。该混合范数界使迟滞输出在 \(L^qL^r\) 中收敛；前面的逐时能量界又使它们一致收敛到已经定义的 \(\mathcal R_I F\)，收敛发生在 \(C_b(\mathbb R;L^2)\) 中。两个极限在时空分布配对中一致，故混合范数极限也是这个连续代表。这同时证明了延拓的唯一性与\eqref{b2:eq:E-retarded-bound}。

对有限起点，由\eqref{b2:eq:E-finite-start-duhamel}，第二项是一条以 \(\mathcal R_I F(a)\in L^2\) 为初值的自由轨道。齐次估计与酉性给出
\[
 \begin{aligned}
 \|\mathcal D_aF\|_{L^\infty L^2}
 +\|\mathcal D_aF\|_{L^qL^r}
 &\leq C\left(\|F\|_{L^{q'}L^{r'}(I)}
                       +\|\mathcal R_I F(a)\|_2\right)\\
 &\leq C'\|F\|_{L^{q'}L^{r'}(I)}.
 \end{aligned}
\]
连续性和起点值直接从定义得到。对光滑输入，两个过去积分相减恰是从 \(a\) 到 \(t\) 的有向积分；一般输入由上述连续延拓保持这个记号与含义。
\end{proof}
同一个定义还给出重启恒等式。对命题中的外力及任意有限 \(a,b,t\)，将 \(B_F(t)-B_F(a)\) 在 \(B_F(b)\) 处分开，便得到
\begin{equation}\label{b2:eq:E-duhamel-restart}
 \mathcal D_aF(t)
 =U(t-b)\mathcal D_aF(b)+\mathcal D_bF(t).
\end{equation}
因此 \(u(t)=U(t-a)f+\mathcal D_aF(t)\) 满足
\[
 u(t)=U(t-b)u(b)+\mathcal D_bF(t).
\]
所有值都取上面构造的连续 \(L^2\) 代表，所以可以在任意有限时刻重启，包括把较晚时刻当作起点向后求解。若只关心某一子区间，积分只使用起点与观察时刻之间的外力，区间外如何延拓不会改变该段的输出。

\ProseParagraphBreak

这里直接证明的是同一个可容许对的输入输出估计。迟滞截断之所以保留界，是因为 \(0<\alpha<1\) 时可以用同一个 HLS 正核控制其绝对值；不能把这个步骤推广为“任意全时间算子界都自动保持迟滞截断”，也不能据此宣布任意两对或所有端点组合成立。

\ProseParagraphBreak

应用中还会遇到时间可积、空间平方可积的误差。它与上面的对偶混合空间没有普遍包含关系，但可以直接使用酉性和齐次估计处理。下面的结果也说明，在两种外力空间同时适用时，Duhamel 记号代表同一个对象。

\Needspace{6\baselineskip}
\begin{proposition}[时间可积的平方可积外力]\label{b2:prop:E-l1l2-forcing}
设整数 \(d\geq1\)，\(U(t)\) 为符号 \(e^{-4\pi^2\mathrm it|\xi|^2}\) 的 Schrödinger 酉群。取 \(2<q<\infty\)、\(2\leq r<\infty\) 且 \(2/q+d/r=d/2\)。设 \(I\subset\mathbb R\) 为区间，\(F:I\to L^2(\mathbb R^d;\mathbb C)\) 强可测且
\(\int_I\|F(s)\|_2\,\mathrm ds<\infty\)，并将 \(F\) 在 \(I\) 外补零为 \(\widetilde F\)。则
\[
 \mathcal R_I F(t)=\int_{s<t}U(t-s)\widetilde F(s)\,\mathrm ds,
 \quad
 \mathcal D_aF(t)=\int_a^tU(t-s)\widetilde F(s)\,\mathrm ds
\]
都可直接按 \(L^2\) 值强积分定义，其中 \(a\in\mathbb R\) 有限，第二个积分按方向取符号。两者均属于 \(C_b(\mathbb R;L^2)\cap L^q(\mathbb R;L^r)\)，且
\begin{equation}\label{b2:eq:E-l1l2-forcing}
 \begin{aligned}
 \|\mathcal R_I F\|_{L^\infty L^2}
 +\|\mathcal R_I F\|_{L^qL^r}
 &\leq C\|F\|_{L^1(I;L^2)},\\
 \|\mathcal D_aF\|_{L^\infty L^2}
 +\|\mathcal D_aF\|_{L^qL^r}
 &\leq C\|F\|_{L^1(I;L^2)}.
 \end{aligned}
\end{equation}
常数与 \(I,a\) 无关。在 \(L^1_tL^2_x\) 与 \(L^{q'}_tL^{r'}_x\) 的交空间上，这些定义与\cref{b2:prop:E-retarded}的连续延拓一致。
\end{proposition}
\begin{proof}
酉性使积分中每个向量的范数等于 \(\|\widetilde F(s)\|_2\)，所以强积分绝对收敛，且两个输出的 \(L^\infty L^2\) 范数各自至多为 \(\|F\|_{L^1L^2}\)。写成
\[
 \mathcal R_I F(t)
 =U(t)\int_{-\infty}^tU(-s)\widetilde F(s)\,\mathrm ds,
\]
内部积分随上限在 \(L^2\) 中连续；酉群强连续性再给出输出连续。有限起点完全相同。

先对光滑紧支外力作混合范数估计。Minkowski 不等式及齐次 Strichartz 界给出
\[
 \begin{aligned}
 \|\mathcal R_I F\|_{L^q_tL^r_x}
 &\leq\int_I
       \|\mathbf1_{\{t>s\}}\cdot U(t-s)F(s)
                              \|_{L^q_tL^r_x}\,\mathrm ds\\
 &\leq C\int_I\|F(s)\|_2\,\mathrm ds.
 \end{aligned}
\]
时间截断只缩小自由轨道的混合范数，而平移 \(t\mapsto t-s\) 不改变全时间齐次常数。由紧支光滑函数在 \(L^1(\mathbb R;L^2)\) 中的稠密性，这个估计定义一般输入的 \(L^qL^r\) 极限；酉性下的积分误差同时给出一致 \(L^2\) 收敛，因此两个极限正是同一强积分输出。

有限起点的积分等于
\(\mathcal R_I F(t)-U(t-a)\mathcal R_I F(a)\)。第二项的混合范数由齐次估计控制，能量范数由酉性控制，得到相同类型的统一界。最后，若 \(F\) 同时属于两种输入空间，先作时空截断，再用光滑化，可以在两种范数中同时逼近它。光滑输入上的定义一致，而两侧各有已证连续界，故交空间上的极限也一致。
\end{proof}
因此可处理写成 \(F=F_1+F_2\) 的外力，其中
\[
 F_1\in L^{q'}_tL^{r'}_x(I),\quad
 F_2\in L^1_tL^2_x(I).
\]
定义 \(\mathcal D_aF=\mathcal D_aF_1+\mathcal D_aF_2\)，交空间上的相容性保证它与分解无关，并有
\[
 \|\mathcal D_aF\|_{C_bL^2\cap L^qL^r}
 \leq C\left(\|F_1\|_{L^{q'}L^{r'}(I)}
                  +\|F_2\|_{L^1L^2(I)}\right),
\]
交空间范数取两项范数之和。后面的非线性构造把主要非线性放入第一类空间；平方可积的时间误差则可由第二类空间控制。两种输入都给出连续的 \(L^2\) 输出，但只有第二种定义预先具有 \(L^2\) 值绝对积分的意义。

\section{波方程的频带衰减与非端点 Strichartz 估计}

\label{b2:sec:app-E-3}

波方程的频率相位与 Schrödinger 方程不同：它对 \(|\xi|\) 是一次齐次的，
沿径向没有二次曲率。因此先前直接写全频率有界函数核的方法不能照搬。
先把输入限制在一个环带，角向曲率便提供一致衰减；
随后把每个频带的时空估计按平方和合成，才能恢复恰当的 Sobolev 初值范数。

\ProseParagraphBreak
以下 \(d\ge2\)，令 \(|D|=\sqrt{-\Delta}\)，其 Fourier 符号为 \(2\pi|\xi|\)。
对 \(f\in L^2(\mathbb R^d;\mathbb C)\)、\(t\in\mathbb R\)，定义半波群
\[
 V(t)f\coloneqq
 \mathcal F_2^{-1}\left(e^{2\pi\mathrm it|\xi|}\widehat f(\xi)\right),
 \quad V(t)=e^{\mathrm it|D|}.
\]
乘子模为一，所以 \(V(t)\) 在 \(L^2\) 上酉，且
\(V(t)V(s)=V(t+s)\)、\(V(t)^*=V(-t)\)。
这里的速度尺度是 \(x\) 与 \(t\) 同时缩放，
不同于 Schrödinger 演化中 \(t\) 按空间尺度的平方缩放。

\Needspace{6\baselineskip}
\begin{lemma}[环带半波核]\label{b2:lem:E-half-wave-kernel}
设整数 \(d\ge2\)，\(V(t)\) 为 \(\mathbb R^d\) 上符号
\(e^{2\pi\mathrm it|\xi|}\) 的半波群。
取实偶函数 \(\psi\in C_c^\infty(\mathbb R^d)\)，满足
\(\operatorname{supp}\psi\subset\{\,\xi:1/4\le|\xi|\le4\,\}\)；
对 \(j\in\mathbb Z\)，令 \(P_j\) 为符号 \(\psi(2^{-j}\xi)\) 的 Fourier 乘子。
则对所有 \(t\in\mathbb R\)，有
\[
 \|V(t)P_j\|_{L^1(\mathbb R^d)\to L^\infty(\mathbb R^d)}
 \le C_{d,\psi}2^{jd}(1+2^j|t|)^{-(d-1)/2}.
\]
常数与 \(j,t\) 无关；将截断函数 \(\psi\) 换为 \(\psi^2\) 后结论仍成立。
\end{lemma}
\begin{proof}
紧支撑平滑乘子给出普通卷积核。换元 \(\xi=2^j\eta\) 得
\[
 K_j(t,x)=2^{jd}K_0(2^jt,2^jx),
 \quad
 K_0(t,x)=\int_{\mathbb R^d}
 e^{2\pi\mathrm i(x\cdot\xi+t|\xi|)}\cdot\psi(\xi)\,\mathrm d\xi.
\]
所以只需证明
\(\sup_x|K_0(t,x)|\le C(1+|t|)^{-(d-1)/2}\)。
若 \(|t|\le1\)，直接以 \(\|\psi\|_1\) 控制积分即可。

设 \(|t|>1\)。先考虑 \(|x|<|t|/2\) 或 \(|x|>2|t|\)。
相位的梯度为
\[
 \nabla_\xi(x\cdot\xi+t|\xi|)
 =x+t\frac{\xi}{|\xi|},
\]
其模至少为 \(\bigl||x|-|t|\bigr|\)，在这两个区域又至少为
\((|x|+|t|)/3\)。
记 \(L=|x|+|t|\)，把相位写为 \(L\phi(\xi)\)。
在固定环带上，\(\phi\) 的所需有限阶导数一致有界，
且 \(|\nabla\phi|\ge1/3\)。
因此沿 \(\nabla\phi\) 方向使用分部积分算子
\[
 \frac{\nabla\phi\cdot\nabla_\xi}
      {2\pi\mathrm iL|\nabla\phi|^2}
\]
反复积分，便得到任意固定整数 \(N\) 所对应的界
\(|K_0(t,x)|\le C_{N,d,\psi}L^{-N}\)。
振幅紧支撑，没有边界项；导数控制与梯度下界说明常数不依赖 \(x,t\)。

剩下 \(|t|/2\le|x|\le2|t|\)。写 \(\xi=\rho\omega\)，其中
\(1/4\le\rho\le4\)、\(\omega\in S^{d-1}\)，得到
\[
 K_0(t,x)=\int_{1/4}^4 e^{2\pi\mathrm it\rho}\rho^{d-1}
 \left(\int_{S^{d-1}}e^{2\pi\mathrm i\rho x\cdot\omega}
             \cdot\psi(\rho\omega)\,\mathrm d\omega\right)\mathrm d\rho.
\]
须对内层角积分取得同时关于 \(\rho\) 和方向 \(x/|x|\) 的一致界。
对每个非零 \(x\)，任选正交变换 \(O\) 使 \(Oe_d=x/|x|\)，再令
\(\omega=O\theta\)。角向相位成为 \(\rho|x|\theta_d\)，
振幅成为 \(\psi(\rho O\theta)\)。
不需要让 \(O\) 随 \(x/|x|\) 连续选择：
正交矩阵的系数有界，\(\rho\) 位于固定紧区间，
故这些振幅在球面固定坐标中的有限阶导数具有同一上界。

函数 \(\theta\mapsto\theta_d\) 在球面上只有南北两极驻点。
取与 \(O,\rho,x\) 无关的固定光滑分割，
在两个极点附近分别用坐标
\(\theta=(z,\pm\sqrt{1-|z|^2})\)。
相位的 Hesse 矩阵在 \(z=0\) 为 \(\mp I_{d-1}\)，
因而行列式绝对值为一；坐标邻域、面积因子及其导数也都固定。
当参数 \(\rho|x|\ge1\) 时，
\cref{b2:thm:21-stationary-phase}及其后说明的有限阶导数控制给出
\(C(\rho|x|)^{-(d-1)/2}\) 的界。
其余球面片上切向梯度有统一正下界，可用非驻相分部积分取得更快衰减。
若 \(\rho|x|<1\)，则在当前区域中 \(|x|\) 仍有正下界且有统一上界，
绝对积分也可纳入同一个 \(C|x|^{-(d-1)/2}\) 界。
由 \(\rho\in[1/4,4]\)，最终得到
\[
 \left|\int_{S^{d-1}}e^{2\pi\mathrm i\rho x\cdot\omega}
                    \cdot\psi(\rho\omega)\,\mathrm d\omega\right|
 \le C_{d,\psi}|x|^{-(d-1)/2}.
\]
对 \(\rho\) 作绝对积分，再用 \(|x|\ge|t|/2\)，便得所需时间衰减。
缩放回第 \(j\) 个环带，并以卷积核的上确界控制
\(L^1\to L^\infty\) 范数，完成证明。
\end{proof}

这个估计保留了两个尺度。当 \(2^j|t|\lesssim1\) 时，核的大小至多是频带体积
\(2^{jd}\)；当 \(2^j|t|\gg1\) 时，角向振荡才产生衰减。
在光锥 \(|x|\approx|t|\) 附近，径向相位仍可能相干，
所以衰减指数来自 \(d-1\) 个角向曲率方向。
Schrödinger 的二次相位则有 \(d\) 个非退化方向；
这一区别随后表现为不同的时间空间指数关系。

\ProseParagraphBreak
波可容许关系通常写成
\[
 \frac1q+\frac{d-1}{2r}\le\frac{d-1}{4},
 \quad q,r\ge2,
\]
并另行处理维数相关的端点。
本节完整证明等号成立且 \(q,r\) 都为有限非端点指数的范围。
先计算初值应具有多少阶正则性。
对频率远离零的光滑 \(f\)，令
\(f_\lambda(x)=\lambda^{d/2-s}f(\lambda x)\)、\(\lambda>0\)。
则其齐次 \(s\) 阶范数不变，而半波缩放给出
\[
 V(t)f_\lambda(x)
 =\lambda^{d/2-s}(V(\lambda t)f)(\lambda x),
 \quad
 \|V(t)f_\lambda\|_{L^q_tL^r_x}
 =\lambda^{d/2-s-d/r-1/q}\|V(t)f\|_{L^q_tL^r_x}.
\]
因此全尺度估计所对应的阶数是
\[
 s=d\left(\frac12-\frac1r\right)-\frac1q.
\]
速度在波解公式中还要经过一次 \(|D|^{-1}\)，所以它应放在低一阶的空间中。
下面把这一低阶空间连同低频含义一起明确下来。

\Needspace{6\baselineskip}
\begin{definition}[本节的齐次 Sobolev 数据空间]\label{b2:def:E-homogeneous-data}
设整数 \(d\ge1\)、实数 \(\sigma<d/2\)，记
\[
 \mathcal D_0=
 \{\,\mathcal F^{-1}h:h\in C_c^\infty(\mathbb R^d\setminus\{0\};\mathbb C)\,\}.
\]
在 \(\mathcal D_0\) 上置
\[
 \|f\|_{\dot H^\sigma}^2
 \coloneqq\int_{\mathbb R^d}
       (2\pi|\xi|)^{2\sigma}|\widehat f(\xi)|^2\,\mathrm d\xi.
\]
其完备化记为 \(\dot H^\sigma(\mathbb R^d)\)，称为本节所用的
\term[qici-sobolev-kongjian]{齐次 Sobolev 空间}[homogeneous Sobolev space]。
\symidx{\(\dot H^\sigma\)：齐次 Sobolev 空间}
\end{definition}
这一完备化有具体的温和分布实现。
加权平方可积函数可先截断在远离零点的有界环带，再用光滑函数逼近，
因此其 Fourier 侧恰是满足
\((2\pi|\xi|)^\sigma h\in L^2\) 的可测函数 \(h\)，按几乎处处相等识别。
对任意 \(\phi\in\mathcal S(\mathbb R^d)\)，加权 Cauchy--Schwarz 给出
\[
 \left|\int_{\mathbb R^d}h(\xi)\cdot\phi(\xi)\,\mathrm d\xi\right|
 \le\|(2\pi|\xi|)^\sigma h\|_2
       \cdot\|(2\pi|\xi|)^{-\sigma}\phi\|_2.
\]
第二个因子在零点附近有限，正因为 \(2\sigma<d\)；
在无穷远处则由 Schwartz 衰减控制，并可由有限个 Schwartz 半范数估计。
所以 \(h\) 定义温和分布，完备化中的极限也唯一给出其逆 Fourier 变换。
这里没有额外任意选择的零频多项式。
不等式同时说明该完备化连续嵌入 \(\mathcal S'\)。

\ProseParagraphBreak
负阶情形也使用同一定义，并不要求输入先属于 \(L^2\)。
例如 \(\dot H^{-1/2}\) 的条件是
\((2\pi|\xi|)^{-1/2}\widehat g\in L^2\)，
它允许某些没有通常平方可积性的高频数据，
同时仍对低频施加相应的加权限制。
乘子 \(|D|^{-1}\) 在 \(\mathcal D_0\) 上满足
\[
 \||D|^{-1}g\|_{\dot H^s}=\|g\|_{\dot H^{s-1}}.
\]
因此只要 \(s<d/2\)，它就唯一延拓为
\(\dot H^{s-1}\to\dot H^s\) 的等距映射。
这是波初速度公式中逆导数的定义，不是在任意分布上未经说明地除以 \(|\xi|\)。

\Needspace{6\baselineskip}
\begin{theorem}[全频率边界非端点半波估计]\label{b2:thm:E-wave-nonendpoint}
设整数 \(d\ge2\)，\(2<q<\infty\)、\(2<r<\infty\)，且
\[
 \frac2q+(d-1)\left(\frac1r-\frac12\right)=0,
 \quad
 s=d\left(\frac12-\frac1r\right)-\frac1q.
\]
令 \(V(t)\) 为 Fourier 符号 \(e^{2\pi\mathrm it|\xi|}\) 的半波算子，
\(\dot H^\sigma(\mathbb R^d)\) 按\cref{b2:def:E-homogeneous-data}的完备化定义。
此时 \(0<s<1\) 且 \(s<d/2\)，对每个 \(f\in\dot H^s(\mathbb R^d)\)，有
\[
 \|V(t)f\|_{L^q(\mathbb R;L^r(\mathbb R^d))}
 \le C_{d,q,r}\|f\|_{\dot H^s}.
\]
若另有 \(g\in\dot H^{s-1}(\mathbb R^d)\)，则由半波公式定义的波解
\[
 u(t)=\frac12\bigl(V(t)+V(-t)\bigr)f
 +\frac1{2\mathrm i}\bigl(V(t)-V(-t)\bigr)|D|^{-1}g
\]
满足 \(u_{tt}-\Delta u=0\)，
\(u\in C(\mathbb R;\dot H^s)\)、\(u_t\in C(\mathbb R;\dot H^{s-1})\)，
并有 \(u(0)=f\)、\(u_t(0)=g\) 及估计
\[
 \|u\|_{L^q(\mathbb R;L^r(\mathbb R^d))}
 \le C_{d,q,r}
       \left(\|f\|_{\dot H^s}+\|g\|_{\dot H^{s-1}}\right).
\]
方程按时空分布理解，初值按上述两个空间的连续性理解。
\end{theorem}
\begin{proof}
记 \(a=1/2-1/r\)。指数关系给出
\[
 \alpha=(d-1)a=\frac2q\in(0,1),
 \quad
 s=\frac{d+1}{2}a.
\]
又因 \(a<1/2\) 且 \(a<1/(d-1)\)，有
\[
 0<s<\min\left(\frac{d+1}{4},\frac{d+1}{2(d-1)}\right).
\]
当 \(d=2\) 时右端为 \(3/4\)，当 \(d=3\) 时为1，
当 \(d\ge4\) 时小于1；故 \(0<s<1\) 且 \(s<d/2\)，
两个数据空间均在前面定义的范围内。

先作单个频带的估计。取核引理中的实偶截断 \(\psi\)，
令 \(T_jh(t)=V(t)P_jh\)。
在紧支撑光滑时空测试函数上，
\[
 T_j^*F=\int_{\mathbb R}P_jV(-\tau)F(\tau)\,\mathrm d\tau,
 \quad
 T_jT_j^*F(t)=
 \int_{\mathbb R}V(t-\tau)P_j^2F(\tau)\,\mathrm d\tau.
\]
实值符号保证 \(P_j^*=P_j\)，而所有 Fourier 乘子彼此交换。
核引理应用于 \(\psi^2\) 给出 \(L^1\to L^\infty\) 界，
Plancherel 则给出与 \(j,t\) 无关的 \(L^2\to L^2\) 界。
两者插值的色散权重为 \(1-2/r=2a\)，所以
\[
 \begin{aligned}
 \|V(t)P_j^2\|_{L^{r'}\to L^r}
 &\le C2^{2jda}(1+2^j|t|)^{-\alpha}\\
 &\le C2^{j(2da-\alpha)}|t|^{-\alpha}
 =C2^{2js}|t|^{-\alpha},
 \quad t\ne0.
 \end{aligned}
\]
这里 \(2da-\alpha=2da-2/q=2s\)，
正是频率代价与尺度计算相接的位置。
Minkowski 不等式及一维 Hardy--Littlewood--Sobolev 定理给出
\[
 \|T_jT_j^*F\|_{L^q_tL^r_x}
 \le C2^{2js}\|F\|_{L^{q'}_tL^{r'}_x}.
\]
所用分数积分阶数为 \(1-\alpha\in(0,1)\)，
且 \(1/q=1/q'-(1-\alpha)\)，故正处于强型适用范围。
再对 \(F\) 配对，有
\[
 \begin{aligned}
 \|T_j^*F\|_2^2
 &=\langle T_jT_j^*F,F\rangle\\
 &\le C2^{2js}\|F\|_{L^{q'}_tL^{r'}_x}^2.
 \end{aligned}
\]
开平方，再作对偶，得到
\[
 \|V(t)P_jh\|_{L^q_tL^r_x}\le C2^{js}\|h\|_2.
\]
不能把 \(TT^*\) 的常数 \(2^{2js}\) 原样当成 \(T_j\) 的常数。
上述恒等式先对测试函数使用强积分和 \(L^2\) 配对，
再由所获界及有限指数混合空间的对偶性延拓；
延拓的输出与原来的半波在分布配对中一致。

现在合成所有频带。选取\chapref{b2:ch:19}的实径向分割
\(\varphi=\chi-\chi(2\cdot)\)，
\(\Delta_jf=\mathcal F^{-1}(\varphi(2^{-j}\xi)\widehat f)\)。
其支撑位于固定倍数的第 \(j\) 个环带，且在 \(\xi\ne0\) 上
\(\sum_j\varphi(2^{-j}\xi)=1\)，各点只有有限个非零项。
取上面的 \(\psi\) 在 \(\operatorname{supp}\varphi\) 上等于一，
则 \(P_j\Delta_j=\Delta_j\)。
由频带估计，对 \(u_j(t)=V(t)\Delta_jf\) 有
\[
 \|u_j\|_{L^q_tL^r_x}\le C2^{js}\|\Delta_jf\|_2.
\]
先考虑 \(f\in\mathcal D_0\) 和有限指标集 \(J\)。
\cref{b2:thm:19-square-function,b2:prop:19-square-synthesis}的环带合成界给出
\[
 \left\|\sum_{j\in J}u_j(t)\right\|_r
 \le C_r\left\|\left(\sum_{j\in J}|u_j(t)|^2\right)^{1/2}\right\|_r.
\]
这里恢复其关键理由：在与 \(h\in L^{r'}\) 配对时，
可在 \(h\) 上插入对第 \(j\) 环带等于一的扩大投影 \(\widetilde\Delta_j\)；
逐点序列 Cauchy--Schwarz 和对偶平方函数界给出
\[
 \left|\sum_{j\in J}\int u_j(t,x)\cdot\overline{h(x)}\,\mathrm dx\right|
 \le C_r
 \left\|\left(\sum_{j\in J}|u_j(t)|^2\right)^{1/2}\right\|_r
 \cdot\|h\|_{r'}.
\]
在 \(L^{r'}\) 单位球上取上确界即得合成不等式。
因为 \(q/2,r/2>1\)，空间和时间中的 Minkowski 不等式又给出
\[
 \begin{aligned}
 \left\|\left(\sum_{j\in J}|u_j|^2\right)^{1/2}\right\|_{L^q_tL^r_x}^2
 &=\left\|\sum_{j\in J}|u_j|^2\right\|_{L^{q/2}_tL^{r/2}_x}\\
 &\le\sum_{j\in J}\||u_j|^2\|_{L^{q/2}_tL^{r/2}_x}
 =\sum_{j\in J}\|u_j\|_{L^q_tL^r_x}^2.
 \end{aligned}
\]
于是
\[
 \left\|\sum_{j\in J}V(t)\Delta_jf\right\|_{L^q_tL^r_x}
 \le C\left(\sum_{j\in J}2^{2js}\|\Delta_jf\|_2^2\right)^{1/2}.
\]
在第 \(j\) 环带上 \(2^{js}\) 与 \((2\pi|\xi|)^s\) 可比，
Plancherel 及有限重叠因而给出
\[
 \begin{aligned}
 \sum_j2^{2js}\|\Delta_jf\|_2^2
 &=\int\left(\sum_j2^{2js}|\varphi(2^{-j}\xi)|^2\right)
                    |\widehat f(\xi)|^2\,\mathrm d\xi\\
 &\le C\|f\|_{\dot H^s}^2.
 \end{aligned}
\]
对 \(\mathcal D_0\) 输入只涉及有限个非零频带，
故求和重构 \(V(t)f\)，并取得定理的半波界。

对一般 \(f\in\dot H^s\)，先由 Fourier 侧的单位模乘子定义
\(V(t)f\in\dot H^s\)。支配收敛说明轨道在该空间连续。
取 \(f_n\in\mathcal D_0\) 趋于 \(f\)，已证估计使
\(V(t)f_n\) 在 \(L^q_tL^r_x\) 中 Cauchy；
同时它们在每个有限时间区间上一致趋于 \(V(t)f\) 于 \(\dot H^s\)。
前面证明的 \(\dot H^s\to\mathcal S'\) 连续嵌入将两个极限识别为同一时空分布，
因而得到一般输入的估计。
也可以直接对任意有限尾指标集应用合成界：
上式平方和的尾部趋零，所以全频率级数在 \(L^q_tL^r_x\) 中无条件收敛。

最后令 \(h=|D|^{-1}g\in\dot H^s\)。
波解公式只含 \(f,h\) 的四个半波，时间反向不改变全线混合范数，
故由三角不等式取得所述估计。
该公式在 \(\mathcal D_0\) 输入上可直接求导，并给出
\(u_{tt}+|D|^2u=0\)、\(u(0)=f\)、\(u_t(0)=g\)。
对一般数据，\(|D|:\dot H^\sigma\to\dot H^{\sigma-1}\) 是等距映射，
而对实数 \(\delta\ne0\)、\(\xi\ne0\) 有
\[
 \left|\frac{e^{2\pi\mathrm i\delta|\xi|}-1}{2\pi\delta|\xi|}\right|\le1.
\]
因此 Fourier 侧的支配收敛使半波求导在低一阶空间中成立。
于是 \(u_t\in C(\mathbb R;\dot H^{s-1})\)，
二阶方程在 \(\dot H^{s-2}\) 及分布意义下成立，
初值也随连续性传到极限。
\end{proof}

输入的频带正交只直接控制平方范数。
输出空间通常是 \(L^r\)、\(r\ne2\)，因此必须用平方函数把输出求和转回可控的平方和。
若改为逐带三角不等式，右端会成为
\(\sum_j2^{js}\|\Delta_jf\|_2\)；
它要求的是更强的频带可和性，不能再称为同一个 \(\dot H^s\) 估计。

\ProseParagraphBreak
三维的 \((q,r)=(4,4)\) 满足边界关系，且 \(s=1/2\)。
因此位移在 \(\dot H^{1/2}\)、速度在 \(\dot H^{-1/2}\) 时，
波解属于 \(L^4(\mathbb R;L^4(\mathbb R^3))\)。
这里负阶速度按加权 Fourier 完备化理解，仍有确定的
\(C_t\dot H^{-1/2}\) 初速度迹。
二维取 \(r=4\) 时则有 \(q=8\)、\(s=3/8\)；
相同空间指数不能保留三维的时间指数。

\ProseParagraphBreak
边界也显示出两种不同障碍。
三维沿上述指数关系令 \(q\downarrow2\) 时，\(r\uparrow\infty\)：
既有时间 HLS 中 \(\alpha\uparrow1\) 的障碍，
又离开有限空间指数的平方函数范围。
四维的候选边界 \((q,r)=(2,6)\)、\(s=5/6\) 则仍有有限的空间指数，
但时间 HLS 障碍依旧存在，所以也不在本节已证定理内。
本节没有把这些证明边界当成统一的失败结论；
系统的其他可容许范围和端点处理见 Keel 与陶哲轩的论文
《Endpoint Strichartz Estimates》\cite{Keel1998Endpoint}。

\section{Schrödinger 能量端点与双线性求和}

\label{b2:sec:app-E-4}

在 Schrödinger 可容许关系 \(2/q+d/r=d/2\) 中令 \(q=2\)，
便得到 \(r=2d/(d-2)\)。当 \(d\ge3\) 时，这是一个有限空间指数；
但是前面的时间 HLS 证明仍然失效，因为色散插值后的时间幂变成了 \(|t-s|^{-1}\)。
失败的是这条正核估计路线。实际的欧氏自由演化仍满足下面的齐次端点结论，
其证明要保留算子配对，并在多个时间尺度之间作双线性求和。

\Needspace{6\baselineskip}
\begin{theorem}[深层结构定理：欧氏 Schrödinger 能量端点]\label{b2:thm:E-endpoint}
设整数 \(d\ge3\)，\(f\in L^2(\mathbb R^d;\mathbb C)\)，并令
\[
 U(t)f=\mathcal F_2^{-1}
       \left(e^{-4\pi^2\mathrm it|\xi|^2}\widehat f(\xi)\right),
 \quad t\in\mathbb R.
\]
则存在只依赖 \(d\) 的常数 \(C_d\)，使
\[
 \|U(t)f\|_{L^2(\mathbb R;L^{2d/(d-2)}(\mathbb R^d))}
 \le C_d\|f\|_2.
\]
二维候选对 \((q,r)=(2,\infty)\) 不属于本定理的结论。
\end{theorem}
\begin{proofsketch}
这里使用 Keel--Tao 抽象定理的齐次端点特例，先明确其所需假设。
设 \(\mathcal H\) 为复 Hilbert 空间，\((Y,\nu)\) 为
\(\sigma\)-有限测度空间，\(W(t):\mathcal H\to L^2(Y,\nu)\) 是可测算子族。
所需的两个一致估计为
\[
 \|W(t)h\|_{L^2(Y)}\le C_E\|h\|_{\mathcal H},
\]
以及对 \(t\ne s\)、\(v\in L^1(Y)\cap L^2(Y)\) 成立的
\[
 \|W(t)W(s)^*v\|_{L^\infty(Y)}
 \le C_D|t-s|^{-\sigma}\|v\|_{L^1(Y)}.
\]
其中时间变量取遍 \(\mathbb R\)，\(\sigma>1\)，
\(C_E,C_D\) 不依赖 \(s,t,h,v\)。
这里的可测性包括所讨论的 \(W(t)h\) 时空函数可测；
\(W(s)^*:L^2(Y)\to\mathcal H\) 由第一条有界性定义。
抽象结果在这些假设下给出
\[
 \|W(t)h\|_{L^2_tL^{r_\sigma}(Y)}
 \le C\|h\|_{\mathcal H},
 \quad r_\sigma=\frac{2\sigma}{\sigma-1}.
\]
这正是尖锐指数关系
\(1/q+\sigma/r=\sigma/2\) 在 \(q=2\) 时的有限空间指数端点。
完整抽象定理还处理其他指数；本处只取上述特例。

对当前问题，取
\(\mathcal H=L^2(\mathbb R^d)\)、\(Y=\mathbb R^d\)、\(\nu=\mathrm dx\)、
\(W(t)=U(t)\)。Plancherel 给出能量界 \(C_E=1\)，
强连续性给出所需可测性。
又因 \(U(s)^*=U(-s)\) 及群律，
\[
 W(t)W(s)^*=U(t-s).
\]
本附录首节证明的 Fresnel 核界给出
\[
 \|U(t-s)v\|_\infty
 \le(4\pi|t-s|)^{-d/2}\|v\|_1,
 \quad t\ne s.
\]
所以可取 \(\sigma=d/2>1\)，从而
\(r_\sigma=2d/(d-2)\)，与定理中的指数一致。
这一应用核对的是全部时间差的色散界，并非只在大时间成立的衰减。

抽象端点定理的证明把 \(TT^*\) 配对按时间间隔作二进分解，
利用邻近空间指数的局部双线性估计，再以双线性插值和尺度序列求和得到临界界。
这些局部界的完整适用邻域和端点插值步骤在此不重证，
也不能用对 \(|t-s|^{-1}\) 直接套用强型 HLS 代替。
正式来源是 Keel 与陶哲轩的
《Endpoint Strichartz Estimates》
\cite[Theorem 1.2, p.~956; Corollary 1.4, p.~958; Sections 4--6]{Keel1998Endpoint}。
其中抽象定理提供能量与色散假设下的结论，
相应推论将它应用到欧氏 Schrödinger 演化。
陶哲轩的《Nonlinear Dispersive Equations: Local and Global Analysis》
\cite[Section 2.3]{Tao2006Nonlinear}亦给出相应背景。
本附录的非线性模型只调用已经完整证明的非端点界，
不以这里略去的端点求和作为隐含前置。
\end{proofsketch}

两个假设在应用中各有作用。能量界使 \(W(s)^*\) 有意义，
并控制没有空间可积性改善时的演化；色散界则改善空间指数，但在对角线 \(t=s\) 附近发散。
抽象定理把这两种控制联合起来。
在欧氏模型中，酉性和群律让 \(W(t)W(s)^*\) 化为单个时间差的演化，
因此前面已经建立的核估计正好满足所需形式；
不能仅凭某个算子族有 \(L^1\to L^\infty\) 衰减就省略对其伴随乘积的核对。

\ProseParagraphBreak
下面具体说明尺度分解为何是实质步骤，而不把这一说明当作端点证明。
置 \(r_*=2d/(d-2)\)，则 \(r_*'=2d/(d+2)\)。
对 \(F,G\in C_c^\infty(\mathbb R\times\mathbb R^d;\mathbb C)\)，考虑双线性配对
\[
 B(F,G)=\int_{s<t}
 \langle U(t-s)F(s),G(t)\rangle_{L^2_x}\,\mathrm ds\,\mathrm dt.
\]
这里内积约定对第一个变量复线性，故复标量下 \(B\) 对第二个变量共轭线性；
双线性求和沿用这一通常的复配对约定。
当前测试函数的时间支撑紧，能量界已保证这一双积分存在；
端点问题是能否仅用
\(L^2_tL^{r_*'}_x\) 范数取得与支撑大小无关的估计。
把上半时间平面按间距分为
\[
 B(F,G)=\sum_{k\in\mathbb Z}B_k(F,G),
 \quad
 B_k(F,G)=
 \int_{2^k\le t-s<2^{k+1}}
 \langle U(t-s)F(s),G(t)\rangle_{L^2_x}\,\mathrm ds\,\mathrm dt.
\]
对测试函数，能量界也使这个分解合法；
困难在于求和后的目标范数控制。

\ProseParagraphBreak
色散与能量插值在 \(r_*\) 处恰给
\[
 \|U(t-s)h\|_{r_*}\le C|t-s|^{-1}\|h\|_{r_*'}.
\]
记 \(f(s)=\|F(s)\|_{r_*'}\)、\(g(t)=\|G(t)\|_{r_*'}\)，
则每一带都受正时间核
\[
 H_k(\tau)=\frac{\mathbf1_{[2^k,2^{k+1})}(\tau)}{\tau},
 \quad
 \|H_k\|_{L^1(\mathbb R)}=\log2
\]
控制。Young 不等式和时间 Cauchy--Schwarz 因而给出
\[
 |B_k(F,G)|
 \le C\|H_k*f\|_2\cdot\|g\|_2
 \le C\|F\|_{L^2_tL^{r_*'}_x}
        \cdot\|G\|_{L^2_tL^{r_*'}_x}.
\]
每个尺度得到的常数都相同；逐带相加便失去可和性。
如果先删去 \(|t-s|<\varepsilon\)，再用正核证明，
仅在 \(\varepsilon<|t-s|<1\) 内就会出现
\(\int_\varepsilon^1\tau^{-1}\,\mathrm d\tau=\log(1/\varepsilon)\)。
因此这种删去不能给出独立于 \(\varepsilon\) 的端点常数。
它揭示的是绝对值估计丢掉了所需结构，
并没有否定原振荡算子本身的有界性。

\ProseParagraphBreak
原典的双线性方法同时允许两个空间指数变化。
为使用两个输入各自的局部时间范数，还把 \(s<t\) 区域分成时间矩形
\(I\times J\)，使两区间的长度与彼此间距同阶。
将长度为 \(2^k\) 的这些矩形贡献合在一起，记为 \(\mathcal B_k\)。
这是对整个双线性配对的另一种分解；
它保留了前面每个时间间距同阶的性质，又使每个输入各自限于一个时间区间。
以 \(\sigma=d/2\) 表示衰减维数，其局部估计在端点附近的适当指数
\(a,b\) 上具有如下尺度形式：
\[
 |\mathcal B_k(F,G)|
 \le C_{a,b}2^{-k\beta(a,b)}
       \|F\|_{L^2_tL^{a'}_x}
       \cdot\|G\|_{L^2_tL^{b'}_x},
 \quad
 \beta(a,b)=\sigma-1-\frac{\sigma}{a}-\frac{\sigma}{b}.
\]
这条局部双线性估计及其可用指数邻域属于上述原典的证明内容，
这里仅用它解释求和机制，并未将它作为本书已经补证的引理。
在 \(a=b=r_*\) 时 \(\beta=0\)；
改变两个指数，可使尺度权重向不同方向衰减。
仅仅在每个尺度先插值回端点，仍会回到 \(\beta=0\) 而没有解决求和。
直接给每个尺度选取较小的邻近指数界也不够：
目标输入只具有临界范数，在整个 \(\mathbb R^d\) 上，
它并不自动属于所有邻近的空间 \(L^{a'}\)、\(L^{b'}\)。
因此还须把输入的大小层次、两个指数的插值和尺度序列的求和放在同一论证中。
原典第4—6节完成的正是这些步骤。

\ProseParagraphBreak
所得双线性控制再返回 Hilbert 空间范数。
对 \(T^*F=\int U(-s)F(s)\,\mathrm ds\)，群律和伴随关系给出
\[
 \|T^*F\|_2^2
 =\int_{\mathbb R^2}
       \langle U(t-s)F(s),F(t)\rangle\,\mathrm ds\,\mathrm dt
 =2\operatorname{Re}B(F,F).
\]
由端点双线性求和得到
\(\|T^*F\|_2^2\le C\|F\|_{L^2_tL^{r_*'}_x}^2\)，
开平方并对偶即为齐次端点估计。
这一收尾与非端点 \(TT^*\) 相同；
新增的深层工作集中在从各个 \(B_k\) 取得整个 \(B\) 的统一界。

\ProseParagraphBreak
结论的范围可以用几个指数核对。
三维得到 \(L^2_tL^6_x\)，四维得到 \(L^2_tL^4_x\)；
二者都没有空间 \(L^\infty\) 端点。
二维对应 \(\sigma=1\)，抽象尖锐指数关系给出的
\((q,r,\sigma)=(2,\infty,1)\) 正是原定理排除的组合，
所以不能把三维公式中的 \(d\) 直接改成2。
一维则不存在满足 Schrödinger 尺度关系的 \(q=2\)、\(r\ge2\) 指数对。
上一节半波的频带衰减维数为 \((d-1)/2\)，
并另带频率代价，也不能将本节的 Schrödinger 端点原样用于波演化。

\ProseParagraphBreak
齐次端点可直接限制到任意时间区间 \(I\)，常数与区间长度无关。
它的对偶还给出
\[
 \left\|\int_{\mathbb R}U(-s)F(s)\,\mathrm ds\right\|_2
 \le C_d\|F\|_{L^2(\mathbb R;L^{r_*'}(\mathbb R^d))},
\]
其中一般输入的积分按测试函数稠密延拓解释。
将齐次算子与这个对偶算子复合，还可得到全时间 \(TT^*\) 的相应界。
但全时间积分与 \(s<t\) 的迟滞积分是不同算子；
不能因为前者有界，就仅凭写入时间示性函数断言全部非齐次端点组合成立。
输出空间 \(r=\infty\)、时间指数 \(q=2\)、两个可容许对同时取端点，
也分别是不同的问题。
本附录随后的一维非线性应用使用已完整证明的非端点同对估计，
其存在性论证不需要在这些未展开的组合之间作选择。

\section{一维三次 Schrödinger 方程的 \texorpdfstring{\(L^2\)}{L²} 全局解}

\label{b2:sec:app-E-5}

现在用已经证明的非端点估计构造非线性演化。
在一维空间中固定 \(\mu\in\{-1,1\}\)，考虑
\[
 \mathrm i\partial_tu+\partial_x^2u=\mu|u|^2u,
 \quad u(0)=u_0\in L^2(\mathbb R;\mathbb C).
\]
符号 \(\mu\) 区分两种非线性，但以下局部收缩和质量守恒对两者都成立。
主要困难是三次项并不把全部 \(L^2\) 映到 \(L^2\)。
我们给解增加一个时间空间混合范数，用线性估计取得这个范数，
再把非线性送入线性非齐次估计允许的输入空间。

\subsection{解空间与三次估计}
\label{b2:subsec:E-cubic-space}

一维可容许对 \((q,r)=(8,4)\) 满足
\(2/8+1/4=1/2\)，处于\cref{b2:thm:E-schrodinger-nonendpoint}的范围。
对有限闭区间 \(I=[a,b]\)、\(a<b\)，记
\[
 \begin{aligned}
 X(I)&\coloneqq C(I;L^2(\mathbb R;\mathbb C))
                   \cap L^8(I;L^4(\mathbb R;\mathbb C)),\\
 \|u\|_{X(I)}&\coloneqq\sup_{t\in I}\|u(t)\|_2
                     +\|u\|_{L^8(I;L^4)},\\
 Y(I)&\coloneqq L^{8/7}(I;L^{4/3}(\mathbb R;\mathbb C)).
 \end{aligned}
\]
交空间中的两个代表须在时间空间上几乎处处相同。
其中连续的 \(L^2\) 代表决定每个时刻的初值，
而 \(L^4\) 值只需在几乎每个时刻存在。
当 \(I=[0,T]\) 时仍写 \(X_T=X(I)\)。

设 \(U(t)\) 为 Fourier 乘子 \(e^{-4\pi^2\mathrm it\xi^2}\) 定义的酉群。
\Cref{b2:prop:E-retarded}把光滑输入上的算子
\[
 \mathcal D_aF(t)\coloneqq\int_a^tU(t-s)F(s)\,\mathrm ds
\]
连续延拓为 \(Y(I)\to X(I)\) 的有界算子。
结合齐次估计，有与 \(a,b\) 无关的常数 \(C_0,C_1\ge1\)，使
\begin{equation}\label{b2:eq:E-cubic-linear-bounds}
 \|U(\,\cdot-a)f\|_{X(I)}\le C_0\|f\|_2,
 \quad \|\mathcal D_aF\|_{X(I)}\le C_1\|F\|_{Y(I)}.
\end{equation}
这里一般的 \(F\) 只属于 \(Y(I)\)，不预先假设
\(F\in L^1(I;L^2)\)。后面所有这类 Duhamel 积分均取上述连续延拓的意义，
其输出仍有连续的 \(L^2\) 代表，并在 \(t=a\) 取零值。

\Needspace{6\baselineskip}
\begin{lemma}[三次项的解空间估计]\label{b2:lem:E-cubic-space}
设 \(I=[a,b]\subset\mathbb R\)、\(a<b\)，
\(X(I)=C(I;L^2(\mathbb R;\mathbb C))\cap L^8(I;L^4(\mathbb R;\mathbb C))\)
取两范数之和，\(Y(I)=L^{8/7}(I;L^{4/3}(\mathbb R;\mathbb C))\)。
则 \(X(I)\) 完备。对 \(u,v\in X(I)\)，令 \(\mathcal N(u)=|u|^2u\)，有
\begin{equation}\label{b2:eq:E-cubic-nonlinear-bounds}
 \begin{aligned}
 \|\mathcal N(u)\|_{Y(I)}
 &\le |I|^{1/2}\cdot\|u\|_{L^8(I;L^4)}^3,\\
 \|\mathcal N(u)-\mathcal N(v)\|_{Y(I)}
 &\le 2|I|^{1/2}
    \left(\|u\|_{L^8(I;L^4)}^2+\|v\|_{L^8(I;L^4)}^2\right)
    \cdot\|u-v\|_{L^8(I;L^4)}.
 \end{aligned}
\end{equation}
\end{lemma}
\begin{proof}
若 \((u_n)\) 在和范数中是 Cauchy 列，两个 Banach 空间分别给出极限
\(u\in C(I;L^2)\) 与 \(v\in L^8(I;L^4)\)。
对任意光滑紧支撑测试函数 \(\phi\)（时间支集位于 \((a,b)\)），有
\[
 \left|\int_I\int_{\mathbb R}(u_n-u)\phi\,\mathrm dx\,\mathrm dt\right|
 \le\|u_n-u\|_{C(I;L^2)}\cdot\|\phi\|_{L^1(I;L^2)}\longrightarrow0.
\]
用 \(L^8L^4\) 与 \(L^{8/7}L^{4/3}\) 配对，同一积分也收敛到以 \(v\) 代替 \(u\) 的值。
因此两个局部可积函数作为分布相同，从而几乎处处相等。
于是连续代表 \(u\) 同时属于两个空间，且 \(u_n\to u\) 于和范数，证明完备性。

对几乎每个 \(t\)，空间 Hölder 给出
\(\|\mathcal N(u(t))\|_{4/3}=\|u(t)\|_4^3\)。
三次方先属于时间空间 \(L^{8/3}(I)\)，再由有限区间上的 Hölder 不等式得到
\[
 \bigl\|\|u(t)\|_4^3\bigr\|_{L^{8/7}(I)}
 \le |I|^{7/8-3/8}\cdot
       \bigl\|\|u(t)\|_4^3\bigr\|_{L^{8/3}(I)}
 =|I|^{1/2}\cdot\|u\|_{L^8(I;L^4)}^3.
\]
对复数 \(z,w\)，把差写成
\(|z|^2(z-w)+(|z|^2-|w|^2)w\)，可得
\[
 \bigl||z|^2z-|w|^2w\bigr|
 \le2(|z|^2+|w|^2)\cdot|z-w|.
\]
先在空间用指数 \(2,4\) 的 Hölder 不等式，
再在时间用 \(4,8\) 把两个因子乘入 \(L^{8/3}(I)\)，
最后使用同一个 \(|I|^{1/2}\) 嵌入因子，就得到差估计。
\end{proof}

对初值 \(f\in L^2\)，所求温和方程是
\begin{equation}\label{b2:eq:E-cubic-mild}
 u(t)=U(t-a)f-\mathrm i\mu\mathcal D_a\mathcal N(u)(t),
 \quad t\in[a,b].
\end{equation}
这里 \(\mathcal N(u)\in Y(I)\) 已由引理保证，故右端是定义良好的 \(X(I)\) 元素。
若 \(a\le c\le t\le b\)，同一解满足重启恒等式
\begin{equation}\label{b2:eq:E-cubic-restart}
 u(t)=U(t-c)u(c)-\mathrm i\mu\int_c^tU(t-s)\mathcal N(u)(s)\,\mathrm ds.
\end{equation}
对光滑外力先把积分在 \(c\) 处分开，再用群律即可验证；
在 \(Y(I)\) 中逼近外力，式\eqref{b2:eq:E-cubic-linear-bounds}保证输出在
\(C(I;L^2)\) 中收敛，因此恒等式也对当前的积分意义成立。
后面从端点继续解时，使用的是这个恒等式。

\subsection{共同寿命与整个解类中的唯一性}
\label{b2:subsec:E-cubic-local}

把式\eqref{b2:eq:E-cubic-mild}的右端看成映射 \(\Phi_f\)。
线性自由轨道的大小由初值质量控制；
三次项的差比线性项多出两个解的范数，
所以要用 \(T^{1/2}\) 抵消这两个因子。
这会给出只依赖初值范数上界的共同时间，而不依赖初值的频率分布。

\Needspace{6\baselineskip}
\begin{theorem}[一维三次模型的局部温和解]\label{b2:thm:E-cubic-local}
固定 \(\mu\in\{-1,1\}\)，令 \(U(t)\) 为
\(L^2(\mathbb R;\mathbb C)\) 上符号为 \(e^{-4\pi^2\mathrm it\xi^2}\) 的酉群。
存在绝对常数 \(c_0,C>0\)，使对任意 \(M>0\)、\(a\in\mathbb R\) 及
\(f\in L^2(\mathbb R;\mathbb C)\)、\(\|f\|_2\le M\)，
令 \(T_M=c_0M^{-4}\)、\(I=[a,a+T_M]\)，则方程
\[
 u(t)=U(t-a)f-\mathrm i\mu\int_a^tU(t-s)(|u|^2u)(s)\,\mathrm ds
\]
在 \(X(I)=C(I;L^2(\mathbb R;\mathbb C))\cap L^8(I;L^4(\mathbb R;\mathbb C))\)
中有唯一解，并有 \(\|u\|_{X(I)}\le CM\)。
积分按\cref{b2:prop:E-retarded}从 \(L^{8/7}(I;L^{4/3})\) 连续延拓的意义理解。
若 \(g\in L^2\)、\(\|g\|_2\le M\)，相同符号和起点对应的解 \(v\) 满足
\[
 \|u-v\|_{X(I)}\le C\|f-g\|_2.
\]
\end{theorem}
\begin{proof}
式\eqref{b2:eq:E-cubic-linear-bounds}与式\eqref{b2:eq:E-cubic-nonlinear-bounds}%
给出一个与区间无关的 \(K\ge1\)，使在长度为 \(T\) 的区间上
\[
 \begin{aligned}
 \|\Phi_f(u)\|_{X(I)}
 &\le C_0M+KT^{1/2}\cdot\|u\|_{X(I)}^3,\\
 \|\Phi_f(u)-\Phi_f(v)\|_{X(I)}
 &\le KT^{1/2}\left(\|u\|_{X(I)}^2+\|v\|_{X(I)}^2\right)
                  \cdot\|u-v\|_{X(I)}.
 \end{aligned}
\]
取 \(R=2C_0M\)，并选绝对常数 \(c_0>0\)，使
\(2KT_M^{1/2}R^2\le1/2\)。
因为 \(T_M^{1/2}R^2=4C_0^2c_0^{1/2}\)，
这个选择同时适用于所有 \(M>0\)。
在闭球 \(\|u\|_{X(I)}\le R\) 中，映射的像范数至多
\(R/2+R/4\)，差估计的系数至多 \(1/2\)。
由完备性和 Banach 不动点定理，闭球中有唯一不动点。
对初值 \(f,g\) 的两个不动点相减，又有
\[
 \|u-v\|_{X(I)}
 \le C_0\|f-g\|_2+\tfrac12\|u-v\|_{X(I)},
\]
吸收后得到 \(2C_0\|f-g\|_2\) 的差界。

还须排除球外的其他 \(X(I)\) 解。
设 \(u,v\in X(I)\) 满足同一温和方程，且在某个左端点 \(c\) 取同值。
在 \(J=[c,d]\subset I\) 上使用重启式，得到
\[
 \|u-v\|_{X(J)}
 \le K|J|^{1/2}
   \left(\|u\|_{L^8(J;L^4)}^2+\|v\|_{L^8(J;L^4)}^2\right)
   \cdot\|u-v\|_{X(J)}.
\]
两个全区间混合范数都有限，故可取统一的小段长，使右侧系数小于一。
把 \(I\) 分成有限个这样的闭区间，第一段从共同初值开始，
差估计使两解相等；连续性给出下一段左端点也相等。
有限次推进证明整个 \(I\) 上唯一。
同一论证适用于任意有限共同存在区间，因此唯一性不限于构造所用的闭球。
\end{proof}

零初值的解为零；需要把它与非零数据比较时，可以把两者一起放进任意正的质量上界 \(M\) 中。
时间平移不改变常数，因而从任何时刻的 \(L^2\) 值重新开始，都有相同的局部结论。
对于一个在 \([0,T_*)\) 上逐紧区间属于 \(X\) 的温和解，若 \(T_*<\infty\)
且存在 \(t_n\uparrow T_*\) 使 \(\|u(t_n)\|_2\le M\)，
则从 \(t_n\) 出发的共同寿命 \(T_M\) 在充分大 \(n\) 时越过 \(T_*\)。
重叠区间上的唯一性使新解接续原解，故 \(T_*\) 不能是最大存在时间。
因此有限最大端点只能伴随 \(L^2\) 范数趋于无穷；下面的质量守恒将排除这种情形。

\subsection{频率投影与带限方程}
\label{b2:subsec:E-cubic-truncation}

仅从 \(u\in X(I)\) 出发，不能直接把微分方程乘以 \(\overline u\) 后分部积分。
先把频率限制在有界区间，使空间导数成为有界算子，并保留投影的自伴性；
后者使近似方程仍然严格守恒质量。

\Needspace{6\baselineskip}
\begin{lemma}[一维区间频率投影]\label{b2:lem:E-frequency-projection}
对 \(N>0\)，令 \(P_N\) 为 \(L^2(\mathbb R;\mathbb C)\) 上 Fourier 乘子
\(\mathbf1_{[-N,N]}\) 对应的正交投影。
对每个 \(1<p<\infty\)，它有兼容的 \(L^p(\mathbb R;\mathbb C)\) 延拓，满足
\[
 \sup_{N>0}\|P_N\|_{L^p\to L^p}\le C_p,
 \quad \|P_Nf-f\|_p\longrightarrow0\quad(N\to\infty)
 \quad(f\in L^p).
\]
若 \(I\) 是有限闭区间、\(1\le q<\infty\)，则对
\(F\in L^q(I;L^p)\) 也有 \(P_NF\to F\) 于该混合范数。
对 \(F\in C(I;L^2)\)，收敛在 \(t\in I\) 上按 \(L^2\) 范数一致。
\end{lemma}
\begin{proof}
由\subsecref{b2:subsec:180502}可知，正频投影为
\(P_+=(I+\mathrm iH)/2\)，其中 \(H\) 是 Hilbert 变换。
其 \(L^p\) 界来自\cref{b2:thm:18-cz-lp}在一维 Hilbert 核上的应用，
并在 \(L^2\cap L^p\) 上与正半轴 Fourier 乘子一致。
调制 \(M_af(x)=e^{2\pi\mathrm iax}f(x)\) 是每个 \(L^p\) 上的等距算子，且
\[
 P_N=M_{-N}P_+M_N-M_NP_+M_{-N}.
\]
右端两个乘子分别是 \(\mathbf1_{\{\xi>-N\}}\) 与
\(\mathbf1_{\{\xi>N\}}\)，端点单点在平方空间中无影响。
所以 \(\|P_N\|_{p\to p}\le2\|P_+\|_{p\to p}\)，常数与 \(N\) 无关。

具有光滑紧支撑 Fourier 变换的 Schwartz 函数在 \(L^p\) 中稠密。
具体地，先用 Schwartz 函数逼近一般 \(L^p\) 输入；
再对其变换乘上 \(\chi(\xi/R)\)，其中 \(\chi\in C_c^\infty\) 在零附近等于一。
由 Schwartz 尾部及截断函数各阶导数的缩放，所得函数随 \(R\to\infty\)
在 Schwartz 拓扑、从而在 \(L^p\) 中趋于原函数。
对于固定的此类函数，充分大 \(N\) 时 \(P_Nf=f\)。
结合统一 \(L^p\) 界，就把强收敛推广到全部 \(L^p\)。

对混合范数输入，几乎每个 \(t\) 上已有空间强收敛，且
\[
 \|(P_N-I)F(t)\|_p^q
 \le(C_p+1)^q\cdot\|F(t)\|_p^q.
\]
右边在时间上可积，支配收敛便给出混合范数结论。
最后，连续映射 \(F:I\to L^2\) 的像为紧集。
在该像上取有限 \(\varepsilon\)-网，
利用 \(\|P_N-I\|_{2\to2}\le2\) 控制网点以外的误差，
再对有限个网点使用强收敛，得到时间一致收敛。
\end{proof}

区间投影使用了一维 Hilbert 变换，不能把任意维数的球形硬截断直接代入这一步。
本节需要的空间指数只有 \(4\) 与 \(4/3\)，两者都在该引理的强型范围内。
此外，\(P_N\) 与 \(U(t)\) 交换；它在 \(Y(I)\) 和 \(X(I)\) 上均有一致界，
所以与 \(\mathcal D_a\) 的交换关系也能从光滑外力通过连续延拓得到。

\Needspace{6\baselineskip}
\begin{proposition}[带限近似的全局解与质量]\label{b2:prop:E-truncated-mass}
设 \(N>0\)、\(\mu\in\{-1,1\}\)、\(f\in L^2(\mathbb R;\mathbb C)\)，
\(P_N\) 为符号 \(\mathbf1_{[-N,N]}\) 的正交投影，
\(\mathcal H_N=P_NL^2(\mathbb R;\mathbb C)\)。
则初值问题
\[
 \mathrm i\partial_tu_N+\partial_x^2u_N
      =\mu P_N(|u_N|^2u_N),
 \quad u_N(0)=P_Nf
\]
有唯一解 \(u_N\in C^1(\mathbb R;\mathcal H_N)\)，方程在 \(L^2\) 中成立。
对全部 \(t\in\mathbb R\)，有
\[
 \|u_N(t)\|_2=\|P_Nf\|_2.
\]
在每个有限闭区间上，它还属于 \(C_tL^2_x\cap L^8_tL^4_x\)。
\end{proposition}
\begin{proof}
闭子空间 \(\mathcal H_N\) 是 Banach 空间，通常仍是无穷维的。
对 \(h\in\mathcal H_N\)，Fourier 反演及 Cauchy--Schwarz 给出
\[
 \|h\|_\infty\le\|\widehat h\|_1
       \le(2N)^{1/2}\cdot\|h\|_2,
 \quad \|\partial_x^2h\|_2\le4\pi^2N^2\cdot\|h\|_2.
\]
由点态三次差估计及这个一致范数界，
\[
 \|P_N(\mathcal N(h)-\mathcal N(k))\|_2
 \le C N\left(\|h\|_2^2+\|k\|_2^2\right)\cdot\|h-k\|_2
 \quad(h,k\in\mathcal H_N).
\]
因此向量场
\(h\mapsto\mathrm i\partial_x^2h-\mathrm i\mu P_N\mathcal N(h)\)
在每个有界球上有界且为 Lipschitz 的，值仍在 \(\mathcal H_N\) 中。
在连续曲线空间中对该向量场的时间积分作收缩：
先以向量场的球上界选时间保证映射留在球内，
再以该球上的 Lipschitz 常数缩短时间保证差的收缩。
这给出一个局部解；因右端连续，解属于 \(C^1_t\mathcal H_N\)。

采用对第一变量线性的内积 \(\langle h,k\rangle=\int h\overline k\)。
带限性使一阶和二阶导数均在 \(L^2\) 中，且 \(u_N\in L^4\)。
由 Plancherel、自伴性及 \(P_Nu_N=u_N\)，
\[
 \begin{aligned}
 \langle\partial_x^2u_N,u_N\rangle
 &=-4\pi^2\int\xi^2|\widehat{u_N}(\xi)|^2\,\mathrm d\xi,\\
 \langle P_N(|u_N|^2u_N),u_N\rangle
 &=\langle|u_N|^2u_N,P_Nu_N\rangle
   =\int|u_N|^4\,\mathrm dx.
 \end{aligned}
\]
这两个数均为实数，故
\[
 \frac{\mathrm d}{\mathrm dt}\|u_N(t)\|_2^2
 =2\operatorname{Re}
   \left(\mathrm i\langle\partial_x^2u_N,u_N\rangle
       -\mathrm i\mu\langle P_N(|u_N|^2u_N),u_N\rangle\right)=0.
\]
对固定 \(N\)，守恒范数给出始终相同的球。
上述局部收缩的时间只依赖 \(N\) 和这个球的半径，
因此可向两个方向重复延拓，不会在有限时间终止。
同时，\(\|h\|_4\le\|h\|_\infty^{1/2}\cdot\|h\|_2^{1/2}
\le(2N)^{1/4}\cdot\|h\|_2\)，
所以解在每个有限时间区间属于 \(L^8_tL^4_x\)。
\end{proof}

这里先固定截断参数得到全局正则解；上述 \(L^4\) 界随 \(N\) 增大，
不能用它直接取极限。
下一步返回 Strichartz 估计，取得不依赖 \(N\) 的短时间控制。

\subsection{质量守恒的极限与全局拼接}
\label{b2:subsec:E-cubic-global}

\Needspace{6\baselineskip}
\begin{theorem}[一维三次模型的质量守恒与全局延拓]
\label{b2:thm:E-cubic-global}
设 \(u_0\in L^2(\mathbb R;\mathbb C)\)、\(\mu\in\{-1,1\}\)，
\(U(t)\) 为符号 \(e^{-4\pi^2\mathrm it\xi^2}\) 的酉群。
温和方程
\[
 u(t)=U(t)u_0-\mathrm i\mu\int_0^tU(t-s)(|u|^2u)(s)\,\mathrm ds
\]
有唯一的全局解
\[
 u\in C(\mathbb R;L^2(\mathbb R;\mathbb C))
       \cap L^8_{\mathrm{loc}}(\mathbb R;L^4(\mathbb R;\mathbb C)),
 \quad \|u(t)\|_2=\|u_0\|_2\quad(t\in\mathbb R).
\]
积分在每个有限时间区间上按\cref{b2:prop:E-retarded}的连续延拓解释，
负时间使用有向积分及时间反演。
这个解延拓\cref{b2:thm:E-cubic-local}给出的局部解；两种符号均适用。
\end{theorem}
\begin{proof}
取 \(M>0\) 使 \(\|u_0\|_2\le M\)，并令 \(u_N\) 为
\cref{b2:prop:E-truncated-mass}的带限解。
在 \(\mathcal H_N\) 中，线性算子 \(\partial_x^2\) 有界，
可对 \(U(-t)u_N(t)\) 在 \(L^2\) 中求导再积分，得到
\[
 u_N(t)=U(t)P_Nu_0-\mathrm i\mu
        \int_0^tU(t-s)P_N\mathcal N(u_N(s))\,\mathrm ds.
\]
此处带限非线性是连续的 \(L^2\) 值函数，故这个积分本身为强积分，
并与它在 \(Y([0,T])\) 上的连续延拓一致。

为把所有截断放到同一个空间中比较，记
\[
 \mathcal N_N(v)\coloneqq P_N\mathcal N(P_Nv),
 \quad
 \Phi_N(v)=U(\,\cdot)P_Nu_0-\mathrm i\mu\mathcal D_0\mathcal N_N(v).
\]
\Cref{b2:lem:E-frequency-projection}在 \(L^4\) 与 \(L^{4/3}\) 上的统一界，
连同三次差估计，给出与 \(N\) 无关的 \(K_*\)，使
\[
 \begin{aligned}
 \|\mathcal N_N(v)\|_{Y(I)}
 &\le K_*|I|^{1/2}\cdot\|v\|_{L^8(I;L^4)}^3,\\
 \|\mathcal N_N(v)-\mathcal N_N(w)\|_{Y(I)}
 &\le K_*|I|^{1/2}
   \left(\|v\|_{L^8(I;L^4)}^2+\|w\|_{L^8(I;L^4)}^2\right)
   \cdot\|v-w\|_{L^8(I;L^4)}.
 \end{aligned}
\]
因此可在共同闭球 \(\|v\|_{X([0,\delta_M])}\le2C_0M\) 上，
选择 \(\delta_M=c_*M^{-4}\le T_M\)，使所有 \(\Phi_N\) 及未截断的 \(\Phi_{u_0}\)
都保持该球，且收缩系数至多 \(1/2\)。
所得 \(\Phi_N\) 的不动点自动带限：\(P_N\) 与 \(U\)、\(\mathcal D_0\) 交换，
对其不动点方程作用 \(P_N\)，右端不变。
另一方面，已构造的 ODE 解 \(u_N\) 属于每个有限区间上的 \(X\)，
并满足同一个积分方程。
用上面的统一差估计重复局部定理中的分段唯一性论证，
便知 ODE 解等于这个不动点。
所以 \(u_N\) 与原方程的局部解 \(u\) 都处于同一个有界球中。

现在固定这个 \(u\)，在 \(I=[0,\delta_M]\) 上有
\(P_Nu\to u\) 于 \(L^8(I;L^4)\)。将误差分成
\[
 \mathcal N_N(u)-\mathcal N(u)
 =P_N\bigl(\mathcal N(P_Nu)-\mathcal N(u)\bigr)+(P_N-I)\mathcal N(u).
\]
第一项由 \(L^{4/3}\) 上的一致投影界、三次差估计及
\(P_Nu\to u\) 控制而趋零；第二项由 \(\mathcal N(u)\in Y(I)\)
及混合范数中的投影强收敛而趋零。
因此
\(\|\mathcal N_N(u)-\mathcal N(u)\|_{Y(I)}\to0\)。
两个不动点方程相减，把依赖 \(u_N-u\) 的一项吸收，得到
\[
 \|u_N-u\|_{X(I)}
 \le2C_0\|P_Nu_0-u_0\|_2
       +2C_1\|\mathcal N_N(u)-\mathcal N(u)\|_{Y(I)}
 \longrightarrow0.
\]
特别地，收敛在整个闭区间上按 \(L^2\) 范数一致。
由近似解的质量守恒及 \(P_Nu_0\to u_0\) 于 \(L^2\)，
对每个 \(t\in[0,\delta_M]\) 有
\[
 \|u(t)\|_2=\lim_{N\to\infty}\|u_N(t)\|_2
          =\lim_{N\to\infty}\|P_Nu_0\|_2=\|u_0\|_2.
\]
这包括右端点，故从 \(\delta_M\) 重新起算时，初值仍满足相同的质量上界。
上述局部构造与守恒证明对任意 \(L^2\) 初值都成立，
所以可以用同一个 \(\delta_M>0\) 逐段继续。
重启恒等式保证拼接后的轨道满足原始温和方程，
局部唯一性保证重叠段相同。
任意有限正时间区间只涉及有限段，故保留 \(L^2\) 连续性与 \(L^8L^4\) 可积性。

负时间可直接由共轭反演构造。
先以 \(\overline{u_0}\) 为初值、同一个 \(\mu\) 构造正向解 \(v\)，
再对 \(t\le0\) 令 \(u(t)=\overline{v(-t)}\)。
Fourier 公式给出 \(\overline{U(s)h}=U(-s)\overline h\)；
对光滑输入变换积分并换元，再用 \(Y\to X\) 的连续延拓，
得到所需的负时间有向温和公式。
它在零时与正向解一致，并保持同一质量。
最后，对任意两个全局解，在每个有限共同区间应用已证唯一性，得到全局唯一性。
\end{proof}

频率截断在这里承担两项不同工作：固定截断后允许严格微分质量，
而截断参数无关的混合范数估计允许取极限。
传递守恒使用了 \(C_tL^2_x\) 中的强收敛；
若只有弱收敛，通常只能得到范数下半连续性，不能从近似质量相等直接取得等式。

\subsection{有限时间稳定性与结论的范围}
\label{b2:subsec:E-cubic-stability}

\Needspace{6\baselineskip}
\begin{corollary}[有限时间上的连续依赖]\label{b2:cor:E-cubic-finite-time-stability}
固定 \(\mu\in\{-1,1\}\)、\(M>0\) 与 \(\tau>0\)。
若 \(f,g\in L^2(\mathbb R;\mathbb C)\)、\(\|f\|_2,\|g\|_2\le M\)，
\(u,v\in C(\mathbb R;L^2)\cap L^8_{\mathrm{loc}}(\mathbb R;L^4)\)
分别为 \(\mathrm i u_t+u_{xx}=\mu|u|^2u\) 从 \(f,g\) 出发的全局温和解，
则存在仅依赖 \(M,\tau\) 的常数 \(C_{M,\tau}\)，使
\[
 \sup_{|t|\le\tau}\|u(t)-v(t)\|_2
 +\|u-v\|_{L^8([ -\tau,\tau];L^4)}
 \le C_{M,\tau}\|f-g\|_2.
\]
\end{corollary}
\begin{proof}
质量守恒使两解在每个重启时刻的 \(L^2\) 范数均不超过 \(M\)。
把 \([0,\tau]\) 分成 \(m\) 段，每段长度不超过局部定理的共同寿命 \(T_M\)。
设局部差界的常数为 \(L\ge1\)。第一段的 \(X\) 差至多
\(L\|f-g\|_2\)，因此下一段的初值差也满足该界。
逐段应用局部定理，得第 \(j\) 段的差至多
\(L^j\|f-g\|_2\)。
全区间连续范数取各段的最大值，混合范数的八次方则是各段八次方之和，
故两者之和至多 \((1+m^{1/8})L^m\cdot\|f-g\|_2\)。
负时间使用共轭反演，作相同估计；两半区间合并即得结论。
\end{proof}

这一定理组的解类是 \(C_tL^2_x\cap L^8_{\mathrm{loc},t}L^4_x\)。
它没有另外证明仅假设 \(C_tL^2_x\) 的任意分布解都自动属于这个类，
也没有把分布方程本身当作温和公式的充分依据。
对于本节构造的解，初值、重启和唯一性均由连续延拓的积分方程确定。

\ProseParagraphBreak
全局存在表示每个有限时间区间都能继续求解；
质量守恒只统一控制 \(L^2\) 范数，
没有给出整个时间轴上的有限 \(L^8_tL^4_x\) 范数或非线性外力的可积尾部。
因此本节不据此推出散射或统一的长期衰减。
一维三次模型的共同寿命来自正的时间小因子 \(T^{1/2}\)；
更改非线性幂次或空间维数时，须重新核对空间乘积的对偶指数和时间嵌入因子，
不能只沿用质量守恒便照搬这个全局论证。

\begin{chaptersummary}
色散界把振荡转成固定时间的空间范数控制，\(TT^*\) 再把这项控制与时间分数积分结合，得到非端点 Strichartz 估计。对偶论证同时给出外力积分的定义和能量界；推迟积分必须独立核对，不能只从全时间算子的有界性删去半个积分域。半波估计先保留环带尺度，再用平方函数把各带的平方和重新合成。端点的正时间核失去可求和性，所引双线性理论保留了取绝对值会丢掉的信息。一维三次模型则展示了另一条组合：统一局部收缩与保质量的正则逼近相接，使平方可积初值得到唯一的全局温和解。
\end{chaptersummary}

\BookExercises

% 自拟：将外力尺度与非线性幂次联系起来，而不只重做齐次缩放。
\Needspace{6\baselineskip}
\begin{exercise}\label{b2:ex:app-E-1}
设整数 \(d\geq1\)，\(u,F\) 为 \(\mathbb R\times\mathbb R^d\) 上满足 \(\mathrm i u_t+\Delta u=F\) 的光滑复值函数。对 \(\lambda>0\)，令
\[
 u_\lambda(t,x)=\lambda^{d/2}u(\lambda^2t,\lambda x).
\]
求相应外力 \(F_\lambda\)，并证明每个 Schrödinger 可容许对 \((q,r)\) 的对偶外力范数 \(L^{q'}_tL^{r'}_x\) 在此缩放下不变，时间均取全实轴。
再对 \(p>1\) 的方程 \(\mathrm i u_t+\Delta u=\mu|u|^{p-1}u\)、\(\mu\in\{-1,1\}\)，求保持方程的振幅指数：令 \(u_\lambda=\lambda^\beta u(\lambda^2t,\lambda x)\)，应取什么 \(\beta\)？由
\(\|f\|_{\dot H^s}=\|(2\pi|\xi|)^s\widehat f\|_2\) 计算临界阶数 \(s_c\)，并判断质量尺度对应的临界幂次。
\end{exercise}
\begin{solution}
时间导数与二阶空间导数各增加因子 \(\lambda^2\)，所以
\[
 F_\lambda(t,x)=\lambda^{d/2+2}F(\lambda^2t,\lambda x).
\]
对一般振幅 \(\lambda^a\)，先在空间换元，再在时间换元，混合范数的缩放因子为 \(\lambda^{a-d/r'-2/q'}\)。本题的指数等于
\[
 \frac d2+2-\frac d{r'}-\frac2{q'}
 =-\frac d2+\frac dr+\frac2q=0.
\]
这里用到 \(1/r'=1-1/r\)、\(1/q'=1-1/q\) 及可容许关系；无穷指数由本性上确界的换元同样得到。若时间只取区间 \(I\)，对应区间必须同时变为 \(\lambda^{-2}I\)，不能固定两个不同尺度的区间却宣称范数恒等。

对幂次方程，左侧缩放为 \(\lambda^{\beta+2}\)，非线性缩放为 \(\lambda^{p\beta}\)。故要求
\[
 \beta+2=p\beta,\quad \beta=\frac2{p-1}.
\]
令 \(f_\lambda(x)=\lambda^\beta f(\lambda x)\)。其 Fourier 变换为 \(\lambda^{\beta-d}\widehat f(\xi/\lambda)\)，从而
\[
 \|f_\lambda\|_{\dot H^s}
 =\lambda^{\beta+s-d/2}\|f\|_{\dot H^s}.
\]
该式先对频率远离零的 Schwartz 函数计算，故无须预设任意齐次空间的实现方式。使幂为零的阶数是
\[
 s_c=\frac d2-\frac2{p-1}.
\]
质量范数对应 \(s=0\)，所以临界幂次为 \(p=1+4/d\)。当 \(1<p<1+4/d\) 时称为质量次临界；此时 \(s_c<0\)，\(L^2\) 所给正则性高于尺度临界阶数。这是尺度判断，本身不证明局部或全局存在。
\end{solution}

% 自拟：窄角频带构造检验波可容许关系的必要性。
\Needspace{6\baselineskip}
\begin{exercise}\label{b2:ex:app-E-2}
设整数 \(d\geq2\)、\(2\leq q,r\leq\infty\)，\(V(t)\) 为 Fourier 符号 \(e^{2\pi\mathrm it|\xi|}\) 的半波算子。假设存在与 \(f\) 无关的常数 \(C\)，使所有 Fourier 支撑在 \(1/2\leq|\xi|\leq3\) 中的 Schwartz 函数都满足
\[
 \|V(t)f\|_{L^q(\mathbb R;L^r(\mathbb R^d))}\leq C\|f\|_2.
\]
取 \(\xi=(\xi_1,\xi')\in\mathbb R\times\mathbb R^{d-1}\)，用支撑在 \(1\leq\xi_1\leq2\)、\(|\xi'|\leq\delta\) 的非负光滑 Fourier 数据，证明必要条件
\[
 \frac1q+\frac{d-1}{2r}\leq\frac{d-1}{4}.
\]
说明这项测试为什么对任何固定阶数的环带 \(\dot H^s\) 右端也适用。
\end{exercise}
\begin{solution}
取非零非负 \(\chi\in C_c^\infty((1,2))\)、\(\eta\in C_c^\infty(\mathbb R^{d-1})\)，令 \(\eta\) 支于单位球，并定义
\[
 \widehat f_\delta(\xi)=\chi(\xi_1)\cdot\eta(\xi'/\delta),
 \quad 0<\delta<\tfrac14.
\]
数据的 Fourier 支撑确实位于题设环带。Plancherel 及横向换元给出
\[
 \|f_\delta\|_2=C_0\delta^{(d-1)/2},\quad C_0>0.
\]
半波相位可写为
\[
 x\cdot\xi+t|\xi|
 =(x_1+t)\xi_1+x'\cdot\xi'
       +t(|\xi|-\xi_1).
\]
在支撑上，\(0\leq|\xi|-\xi_1=|\xi'|^2/(|\xi|+\xi_1)\leq\delta^2/2\)。取充分小且固定的 \(c>0\)，在区域
\[
 0\leq t\leq c\delta^{-2},\quad
 |x_1+t|\leq c,\quad |x'|\leq c\delta^{-1}
\]
上，完整相位 \(2\pi(x\cdot\xi+t|\xi|)\) 的绝对值不超过 \(\pi/3\)。因为振幅非负，积分的实部至少是振幅积分的一半，故
\[
 |V(t)f_\delta(x)|\geq c_1\delta^{d-1}.
\]
每个上述时刻的空间截面体积是常数乘 \(\delta^{-(d-1)}\)，时间长度是常数乘 \(\delta^{-2}\)。于是
\[
 \|V(t)f_\delta\|_{L^qL^r}
 \geq c_2\delta^{d-1-(d-1)/r-2/q}.
\]
无穷指数按对应上确界解释，公式仍成立。与题设上界比较，并令 \(\delta\downarrow0\)，左边的幂不能小于 \((d-1)/2\)。整理即得所求关系。

在固定环带上，\((2\pi|\xi|)^s\) 对任意固定实数 \(s\) 都有与 \(\delta\) 无关的正上下界，所以 \(\|f_\delta\|_{\dot H^s}\) 与 \(\|f_\delta\|_2\) 一致等价。这项构造检测的是角向集中，不能用增加固定环带内的径向 Sobolev 权重消除。
\end{solution}

% 自拟：量化混合范数的次序，避免将时间空间积分随意交换。
\Needspace{6\baselineskip}
\begin{exercise}\label{b2:ex:app-E-3}
设整数 \(d\geq1\)，\(u:\mathbb R\times\mathbb R^d\to\mathbb C\) 可测，且 \(1\leq r\leq q<\infty\)。证明
\[
 \|u\|_{L^q_tL^r_x}\leq\|u\|_{L^r_xL^q_t}.
\]
当 \(q>r\) 时，构造一族有界、有限测度支撑的函数，证明不存在与 \(u\) 无关的常数 \(C\)，使相反方向的估计
\(\|u\|_{L^r_xL^q_t}\leq C\|u\|_{L^q_tL^r_x}\) 对所有这样的函数成立。解释 \(q=r\) 的情形。
\end{exercise}
\begin{solution}
将 \(|u|^r\) 视为关于 \(x\) 积分、取值于时间空间 \(L^{q/r}\) 的非负函数。Minkowski 积分不等式给出
\[
 \begin{aligned}
 \|u\|_{L^q_tL^r_x}^r
 &=\left\|\int_{\mathbb R^d}|u(\,\cdot,x)|^r\,\mathrm dx\right\|_{L^{q/r}_t}\\
 &\leq\int_{\mathbb R^d}
       \bigl\||u(\,\cdot,x)|^r\bigr\|_{L^{q/r}_t}\,\mathrm dx
 =\|u\|_{L^r_xL^q_t}^r.
 \end{aligned}
\]
若右端可能无穷，先对幅值和支撑作有限截断，再由单调收敛得到同一结论。

为否定反向估计，取互不相交、各长一的时间区间 \(I_1,\ldots,I_N\)，以及互不相交、各体积一的空间立方体 \(Q_1,\ldots,Q_N\)，令
\[
 u_N(t,x)=\sum_{k=1}^N\mathbf1_{I_k}(t)\cdot\mathbf1_{Q_k}(x).
\]
固定属于某个 \(I_k\) 的时间时，空间 \(L^r\) 范数为一；固定属于某个 \(Q_k\) 的空间点时，时间 \(L^q\) 范数为一。因此
\[
 \|u_N\|_{L^q_tL^r_x}=N^{1/q},\quad
 \|u_N\|_{L^r_xL^q_t}=N^{1/r}.
\]
比值 \(N^{1/r-1/q}\) 在 \(q>r\) 时无界。若 \(q=r\)，两范数相等，由 Tonelli 定理都等于整个乘积空间的 \(L^q\) 范数。这说明积分次序属于估计的内容，不能把一般 Strichartz 混合范数写成不分次序的双重积分。
\end{solution}

% 自拟：将端点正核的对数下界加强为确切算子范数。
\Needspace{6\baselineskip}
\begin{exercise}\label{b2:ex:app-E-4}
设 \(R>2\)，在实时间轴上定义
\[
 k_R(t)=|t|^{-1}\cdot\mathbf1_{\{1\leq|t|\leq R\}}(t).
\]
证明卷积算子 \(K_Rf=k_R*f\) 在复 \(L^2(\mathbb R)\) 上的算子范数恰为 \(2\log R\)。再求
\(h_\varepsilon(t)=|t|^{-1}\cdot\mathbf1_{\{\varepsilon\leq|t|\leq1\}}(t)\)、\(0<\varepsilon<1\)，所定义卷积算子的范数。说明这些结论反驳的是端点证明中的哪一步，是否已经反驳 Schrödinger 齐次端点估计。
\end{exercise}
\begin{solution}
核非负且 \(\|k_R\|_1=2\log R\)，Young 不等式先给 \(\|K_R\|_{2\to2}\leq2\log R\)。取 \(L>R\) 并令 \(f_L=\mathbf1_{[-L,L]}\)。当 \(|t|\leq L-R\) 时，对核支撑内的每个 \(s\) 都有 \(t-s\in[-L,L]\)，所以
\[
 (k_R*f_L)(t)=2\log R.
\]
故
\[
 \frac{\|K_Rf_L\|_2}{\|f_L\|_2}
 \geq2\log R\sqrt{\frac{L-R}{L}}.
\]
固定 \(R\) 后令 \(L\to\infty\)，得到匹配的下界，因而范数确为 \(2\log R\)。选长区间是为了使远离边界的输入几乎保持常值；不需要常值函数本身属于 \(L^2\)。

对 \(h_\varepsilon\)，其质量为 \(2\log(1/\varepsilon)\)，支撑包含于 \([-1,1]\)。同样用 \(f_L\)、\(L>1\)，在 \(|t|\leq L-1\) 上取得全部核质量，再令 \(L\to\infty\)，得到
\[
 \|h_\varepsilon*\|_{2\to2}=2\log(1/\varepsilon).
\]
两式分别显示长时间与短时间尺度累积的对数代价。若端点 \(TT^*\) 证明把算子配对全数取绝对值后只剩 \(|t-s|^{-1}\)，这些截断不可能给出与截断尺度无关的强 \(L^2\) 界。但正核上界丢掉了振荡与两输入的配对结构；它的失败不推出原演化算子端点估计失败。
\end{solution}

% 自拟：从三次模型推广局部收缩的指数计算，并检验质量临界边界。
\Needspace{6\baselineskip}
\begin{exercise}\label{b2:ex:app-E-5}
设 \(1<p\leq5\)、\(\mu\in\{-1,1\}\)，考虑实线上的复值方程
\[
 \mathrm i u_t+u_{xx}=\mu|u|^{p-1}u,\quad u(0)=u_0\in L^2(\mathbb R).
\]
取 \(r=p+1\)，求使 \((q,r)\) 成为 Schrödinger 可容许对的 \(q\)。在 \([0,T]\) 上计算
\(\||u|^{p-1}u\|_{L^{q'}L^{r'}}\) 由 \(\|u\|_{L^qL^r}^p\) 控制时的时间因子，并证明相应的两输入差估计。当 \(p<5\) 时，说明本附录的收缩法如何给出只依赖初值 \(L^2\) 范数的正存在时间；当 \(p=5\) 时，这一论证的哪项控制消失？
\end{exercise}
\begin{solution}
由 \(2/q+1/(p+1)=1/2\)，得到
\[
 q=\frac{4(p+1)}{p-1},\quad r'=\frac{p+1}{p}.
\]
在所给范围内 \(6\leq q<\infty\)、\(2<r\leq6\)，都位于已经证明的非端点有限指数范围。空间 Hölder 给出
\(\||u(t)|^{p-1}u(t)\|_{r'}=\|u(t)\|_r^p\)。时间上先得到 \(L^{q/p}\)；因为
\[
 \theta\coloneqq\frac1{q'}-\frac pq
 =1-\frac{p+1}{q}=\frac{5-p}{4}\geq0,
\]
有限区间上的 Hölder 不等式给
\[
 \||u|^{p-1}u\|_{L^{q'}L^{r'}}
 \leq T^\theta\|u\|_{L^qL^r}^p.
\]
把复平面视为 \(\mathbb R^2\)，映射 \(z\mapsto|z|^{p-1}z\) 的实导数范数至多 \(p|z|^{p-1}\)。沿 \(w\) 到 \(z\) 的线段积分，便有
\[
 \bigl||z|^{p-1}z-|w|^{p-1}w\bigr|
 \leq C_p\bigl(|z|^{p-1}+|w|^{p-1}\bigr)\cdot|z-w|.
\]
分别对空间与时间应用 Hölder，再用上面的区间嵌入，得到
\[
 \begin{aligned}
 &\bigl\||u|^{p-1}u-|v|^{p-1}v\bigr\|_{L^{q'}L^{r'}}\\
 &\quad\leq C_p T^\theta
   \bigl(\|u\|_{L^qL^r}^{p-1}+\|v\|_{L^qL^r}^{p-1}\bigr)
   \cdot\|u-v\|_{L^qL^r}.
 \end{aligned}
\]
若 \(p<5\)，则 \(\theta>0\)。在 \(C_tL^2\cap L^q_tL^r_x\) 中取半径 \(R=2C_pM\)、\(\|u_0\|_2\leq M\)、\(M>0\)。已有齐次和非齐次界表明，选 \(T\) 使 \(C_pT^\theta R^{p-1}\) 足够小，即得到保球和压缩。可取 \(T=c_p M^{-4(p-1)/(5-p)}\)，其中小常数吸收各线性估计常数。零初值可直接取零解，或置于任意 \(M>0\) 的共同球内。

当 \(p=3\) 时，\((q,r)=(8,4)\)、\(\theta=1/2\)；当 \(p=5\) 时，\((q,r)=(6,6)\)、\(\theta=0\)。五次情形缩短区间不会使这个固定半径球上的系数 \(C_pR^{p-1}\) 变小。这个结论只指出当前保球选择的限制；对给定初值还可利用自由解在短区间的 \(L^6L^6\) 范数趋零，不能据此断言临界方程没有局部解。
\end{solution}

% 自拟：同时排除任意正阶的固定时间正则增益。
\Needspace{6\baselineskip}
\begin{exercise}\label{b2:ex:app-E-6}
设整数 \(d\geq1\)、\(s\geq0\)，定义 \(H^s(\mathbb R^d)\) 的范数为
\[
 \|f\|_{H^s}^2=\int_{\mathbb R^d}
       (1+4\pi^2|\xi|^2)^s|\widehat f(\xi)|^2\,\mathrm d\xi.
\]
构造一个 \(f\in H^s\)，使它不属于任何 \(H^{s+\varepsilon}\)、\(\varepsilon>0\)。证明 \(U(t)=e^{\mathrm it\Delta}\) 对此初值在每个固定实时间也不产生任何正阶的额外 Sobolev 正则性，并解释这与非端点 Strichartz 估计是否矛盾。
\end{exercise}
\begin{solution}
令
\[
 \widehat f(\xi)=\mathbf1_{\{|\xi|\geq e\}}(\xi)
       \cdot|\xi|^{-s-d/2}(\log|\xi|)^{-1}.
\]
因 \(s\geq0\)，该函数平方可积，Plancherel 因而定义唯一的 \(f\in L^2\)。在 \(|\xi|\geq e\) 上，Bessel 权重与 \(|\xi|^{2s}\) 相比只有正的固定常数因子。极坐标积分于是给
\[
 \|f\|_{H^s}^2\asymp\int_e^\infty\frac{\mathrm d\rho}{\rho(\log\rho)^2}<\infty.
\]
对任意 \(\varepsilon>0\)，增加的权重使积分变为
\[
 \int_e^\infty\frac{\rho^{2\varepsilon-1}}{(\log\rho)^2}\,\mathrm d\rho
 =\int_1^\infty\frac{e^{2\varepsilon y}}{y^2}\,\mathrm dy=\infty,
\]
其中最后用 \(y=\log\rho\)。因此同一个输入同时排除了所有正的正则增益。

演化的 Fourier 乘子模恒为一，故每一个上述加权平方积分在每个 \(t\) 都保持原值，包括其无穷性。于是 \(U(t)f\notin H^{s+\varepsilon}\) 对所有 \(t\in\mathbb R\)、\(\varepsilon>0\) 成立。Strichartz 估计控制的是混合可积性，并没有声称输出具有额外空间导数；它与这个例子相容。
\end{solution}

% 自拟：由对偶估计得到附加外力可积性下的散射，区分全局存在与散射。
\Needspace{6\baselineskip}
\begin{exercise}\label{b2:ex:app-E-7}
设整数 \(d\geq1\)，\(U(t)=e^{\mathrm it\Delta}\)，\(2<q<\infty\)、\(2\leq r<\infty\)、\(2/q+d/r=d/2\)。设 \(u_0\in L^2(\mathbb R^d)\)，外力满足 \(F\in L^{q'}([0,\infty);L^{r'}(\mathbb R^d))\)，令
\[
 u(t)=U(t)u_0-\mathrm i\int_0^tU(t-s)F(s)\,\mathrm ds,
\]
积分取本附录已经建立的连续延拓意义。证明存在唯一 \(u_+\in L^2\)，使
\[
 \|u(t)-U(t)u_+\|_2\longrightarrow0\quad(t\to\infty),
\]
并给出由 \(F\) 的尾部范数控制的误差界。作为应用，说明一维三次全局解若另满足 \(u\in L^{24/7}([0,\infty);L^4)\)，便有上述自由渐近状态；质量守恒是否已足以验证这个附加条件？
\end{exercise}
\begin{solution}
由 \(T^*\) 估计，有限 \(R>S\geq0\) 时有
\[
 \left\|\int_S^RU(-s)F(s)\,\mathrm ds\right\|_2
 \leq C\|F\|_{L^{q'}([S,R];L^{r'})}.
\]
这是对偶外力的积分，不预先要求被积函数逐时刻属于 \(L^2\)。有限 \(q'\) 下，右端在 \(S\to\infty\) 时一致趋零，故 \(\int_0^R U(-s)F(s)\,\mathrm ds\) 在 \(L^2\) 中为 Cauchy 族。记其极限为 \(h_\infty\)，定义
\[
 u_+=u_0-\mathrm i h_\infty.
\]
把温和公式乘以有界的 \(U(-t)\)，可先对光滑外力验证，再由连续延拓得到
\[
 U(-t)u(t)-u_+=\mathrm i\int_t^\infty U(-s)F(s)\,\mathrm ds.
\]
酉性与同一尾部估计给出
\[
 \|u(t)-U(t)u_+\|_2
 \leq C\|F\|_{L^{q'}([t,\infty);L^{r'})}\longrightarrow0.
\]
若另一个 \(v_+\) 也满足渐近关系，则酉性和三角不等式使
\(\|u_+-v_+\|_2\leq\|U(t)u_+-u(t)\|_2+\|u(t)-U(t)v_+\|_2\to0\)，所以状态唯一。这种渐近性质称为向前散射到自由解。

一维三次情形取 \((q,r)=(8,4)\)。空间乘积恒等式给
\(\||u(t)|^2u(t)\|_{4/3}=\|u(t)\|_4^3\)，因而
\[
 \||u|^2u\|_{L^{8/7}L^{4/3}}
 =\|u\|_{L^{24/7}L^4}^3.
\]
所给附加条件正好使外力 \(\mu|u|^2u\) 满足前面的全正时间假设。质量守恒只给每个时刻的 \(L^2\) 范数，局部理论只给有限区间的 \(L^8L^4\) 范数；二者都没有给出这个全时间可积性，故此处是条件性结论，不能把它读成全部三次解都散射。
\end{solution}

% 自拟：在温和积分上证明Galilean协变，避免先假设粗糙解可逐点求导。
\Needspace{6\baselineskip}
\begin{exercise}\label{b2:ex:app-E-8}
设 \(\mu\in\{-1,1\}\)，\(u\in C(\mathbb R;L^2(\mathbb R))\cap L^8_{\mathrm{loc}}(\mathbb R;L^4(\mathbb R))\) 是一维方程 \(\mathrm i u_t+u_{xx}=\mu|u|^2u\) 的温和解，初值为 \(u_0\)。对实数 \(v\)，令
\[
 u_v(t,x)=e^{\mathrm i(vx-v^2t)}u(t,x-2vt).
\]
直接从 Fourier 群和温和公式证明它是同号方程、初值 \(e^{\mathrm ivx}u_0\) 对应的解。核验其质量及有限时间混合范数，并判断这个变换是否保持初值的 \(H^1\) 半范数。
\end{exercise}
\begin{solution}
记 \(M_vf(x)=e^{\mathrm ivx}f(x)\)，以及
\(G_v(t)g(x)=e^{\mathrm i(vx-v^2t)}g(x-2vt)\)。在 Schwartz 输入上，调制把 Fourier 频率平移 \(v/(2\pi)\)，而
\[
 -4\pi^2t\left(\xi+\frac v{2\pi}\right)^2
 =-4\pi^2t\xi^2-4\pi tv\xi-tv^2.
\]
在逆变换中换频率变量便得到
\[
 U(t)M_v=G_v(t)U(t).
\]
两侧在 \(L^2\) 上均有界，故由稠密性推广到全部平方可积输入。再用群律，有
\[
 G_v(t)U(t-s)=U(t-s)G_v(s).
\]
对光滑且紧支撑的外力，等式可移入有限时间积分。一般 \(L^{8/7}_tL^{4/3}_x\) 外力以这种数据逼近；\(G_v\) 在每个时间保持空间范数，也保持对应混合范数，已有非齐次连续界遂将恒等式推广到本题的积分意义。

由于单位模乘法和空间平移都与三次非线性相容，
\[
 G_v(s)(|u(s)|^2u(s))=|G_v(s)u(s)|^2G_v(s)u(s).
\]
对原温和公式作用 \(G_v(t)\)，即得
\[
 u_v(t)=U(t)M_vu_0-\mathrm i\mu
       \int_0^tU(t-s)(|u_v|^2u_v)(s)\,\mathrm ds.
\]
对负时间使用同一有向积分恒等式。时间依赖的平移在 \(L^2\) 上强连续，结合 \(u\) 的连续性说明 \(u_v\in C_tL^2\)；空间范数保持还给出任意时间区间 \(I\) 上的
\[
 \|u_v(t)\|_2=\|u(t)\|_2,\quad
 \|u_v\|_{L^q(I;L^r)}=\|u\|_{L^q(I;L^r)}
\]
（双方均可为无穷）。因此它属于本题解类，由唯一性也是变换后初值的全局解。

若 \(u_0\in H^1\)，则
\[
 (M_vu_0)'=e^{\mathrm ivx}(u_0'+\mathrm ivu_0),\quad
 \|(M_vu_0)'\|_2=\|u_0'+\mathrm ivu_0\|_2.
\]
这一般不等于 \(\|u_0'\|_2\)。例如对非零实值 Gaussian 初值，交叉项的实部为零，平方半范数增加 \(v^2\|u_0\|_2^2\)。相位变换保持质量，却会改变导数中的频率中心。
\end{solution}
